\documentclass[review]{elsarticle}
%\documentclass[3p,times,twocolumn,authoryear]{elsarticle}

\usepackage{amssymb}
\usepackage{multirow}
\usepackage{booktabs}
\usepackage[colorlinks,urlcolor=blue]{hyperref}
\usepackage{subfig}
\usepackage{scalefnt}
\usepackage[brazilian]{babel}
\usepackage[utf8]{inputenc}
\usepackage[T1]{fontenc}
\usepackage{units}
\usepackage{url}
\usepackage{amsbsy}
\usepackage{amstext}
\usepackage{graphicx}
\usepackage{rotating}
\usepackage{xcolor}
\usepackage{float}
\usepackage[section]{placeins}
\usepackage{verbatim}
\usepackage{apalike}
\restylefloat{figure}
\restylefloat{table}

% Inserted by pgirardi
\usepackage{amsmath}
\usepackage[normalem]{ulem}
% ponytail: \texttt só — \seqsplit quebra com \_ / babel brazilian
\newcommand{\code}[1]{\texttt{#1}}

% \newcommand{\sugestao}[1]{\textcolor{blue!100!black}{\textbf{#1}}}

\newcommand{\sugestao}[2]{%
	\textcolor{red}{\sout{\textbf{#1}}}%
	\textcolor{blue!100!black}{\textbf{#2}}%
}

%%%%%%%%%%%%%%%

\makeatletter

\journal{Expert Systems with Applications}

\@ifundefined{showcaptionsetup}{}{%
	\PassOptionsToPackage{caption=false}{subfig}}
\usepackage{pdflscape}
\@ifundefined{showcaptionsetup}{}{%
	\PassOptionsToPackage{caption=false}{subfig}}
\makeatother


\begin{document}
	\begin{frontmatter}
		
		\title{Níveis longitudinais hipocampais para sMCI versus pMCI:
			famílias de atributos, quatro coortes e fusão tardia}
		
		\author[ufscar]{Paulo Matheus da Costa Girardi}
		\ead{paulogirardi@estudante.ufscar.br}
		
		\author[ufscar]{Ricardo José Ferrari\corref{cor1}}
		\ead{rferrari@ufscar.br}
		
		\cortext[cor1]{Corresponding author}
		
		\address[ufscar]{Department of Computing, Federal University of São Carlos, Rod.\ Washington Luis, Km 235, 13565-905, São Carlos, SP, Brazil}
		
		\author[]{for the Alzheimer's Disease Neuroimaging Initiative\corref{adni}}
		\cortext[adni]{Data used in preparation of this article were obtained from the Alzheimer's Disease Neuroimaging Initiative (ADNI) database (\url{adni.loni.usc.edu}). Investigators within ADNI contributed to the design and implementation of ADNI and/or provided data but did not participate in analysis or writing of this report. A complete listing of ADNI investigators is available at \url{http://adni.loni.usc.edu/wp-content/uploads/how_to_apply/ADNI_Acknowledgement_List.pdf}.}
		
		\begin{abstract}
			Distinguir sMCI de pMCI no hipocampo com RM T1 isolada é difícil: o sinal é subtil e a literatura mistura janelas, ROIs e validação.
			Neste estudo (ADNI-1/2/GO) pergunta-se se a representação longitudinal \code{t1\_d21\_d32} (Q4: níveis T1, D21 e D32 --- os mesmos descritores em $T_1,T_2,T_3$, \textbf{não} diferenças $T_b-T_a$) agrega discriminação relativamente ao baseline \code{t1\_only}.
			O contraste corre em cinco famílias (volume, forma, textura GLCM Original, primeira ordem, deslocamento normativo) e em quatro coortes (janelas 36/48\,meses $\times$ intervalos nominais 6/12\,meses, tolerância $\pm 2$).
			Protocolo fixo: SVM, pool L1 por frequência (\code{l1\_stable}, $\tau=70\%$, $|\rho|>0.85$ com desempate $|\mathrm{corr}(x,y)|$), CV aninhada no paciente ($5\times 10$); métrica primária AUC de Mann--Whitney sobre a média das probabilidades OOF por sujeito.
			Uniclasse: Q4 agrega volume, textura e primeira ordem; a forma em $t_0$ permanece o teto unimodal e Q4 \textbf{não} o ultrapassa de forma consistente.
			A união tardia de cinco SVMs (média de scores) --- em particular a âncora forma@$t_0$ $\cup$ restante Q4 --- supera esse teto nas quatro células (em \code{48m\_12m}, $n=120$: teto $\approx 0.76$ $\to$ $0.802\pm0.040$).
			Os escores clínicos de $t_0$ (sexo, idade, MMSE, ADAS, FAQ) atingem $0.812\pm0.040$; a RM hipocampal entra como adjunto estrutural.
		\end{abstract}
		
		\begin{keyword}
			Longitudinal MRI \sep hippocampus \sep Mild Cognitive Impairment \sep Radiomics \sep Machine Learning \sep Nested cross-validation \sep Late fusion
		\end{keyword}
		
	\end{frontmatter}
	
	% =============================================================================
	\section{Introdução}
	\label{sec:introducao}
	
	A doença de Alzheimer (AD) e o comprometimento cognitivo leve (MCI) situam-se num contínuo de declínio~\cite{petersen-1999}.
	Identificar quais indivíduos com MCI permanecerão estáveis (sMCI) e quais progredirão para demência (pMCI) interessa a enriquecimento de ensaios, aconselhamento e planeamento de intervenção~\cite{petersen-2009}.
	A RM estrutural T1-ponderada permanece central nesse problema, e a atrofia hipocampal é um dos marcadores macroscópicos mais consolidados~\cite{scheltens-1992}.
	Na fronteira sMCI/pMCI, contudo, a discriminação a partir de uma única região é intrínsecamente difícil: as diferenças de grupo são subtis, a dinâmica de conversão é heterogénea e a alteração estrutural sobrepõe-se ao envelhecimento.
	
	Uma literatura ampla recorreu a representações mais ricas --- morfometria de cérebro inteiro, fusão com escores cognitivos, radiómica e aprendizagem profunda longitudinal~\cite{moradi-2015,ye-2012,zhang-2019,lin-2020,aghajanian-2025}.
	As AUC reportadas podem ser altas, mas a comparabilidade é limitada.
	Coortes, definições de desfecho (p.ex.\ ncMCI/cMCI, conversão MCI$\to$AD num horizonte fixo), escala espacial (cérebro inteiro versus ROI), entrada temporal (só baseline versus visitas repetidas) e desenho de validação diferem.
	Além disso, muitos pipelines aplicam selecção de atributos, ajuste de hiperparâmetros ou limiar de decisão com informação do conjunto completo ou das partições de teste, e nem todos impõem divisão estrita no nível do paciente.
	Sob essas condições, AUC otimistas são difíceis de interpretar e podem não reflectir o desempenho atingível com validação no paciente e ajuste fold-wise.
	
	Este estudo avalia um fluxo auditável de extração de atributos hipocampais e classificação prognóstica na ADNI (ADNI-1/2/GO).
	O hipocampo foi escolhido de propósito como marcador canónico de atrofia e como banco de teste longitudinal, não como cobertura prognóstica de cérebro inteiro.
	A cadeia documentada cobre pré-processamento T1, controlo de qualidade, parcelamento, registro a templates e extração de atributos volumétricos, morfométricos, texturais e de campo de deslocamento no hipocampo.
	O desempenho avalia-se com validação cruzada aninhada no paciente; o modo \code{l1\_stable} (pool L1 por frequência; Sec.~\ref{sec:feature-selection}) restringe-se ao treino de cada dobra externa.
	O classificador canónico é um único SVM, para isolar o contraste entre representações (estática vs longitudinal vs união tardia) de vieses de algoritmo.
	A novidade está na arquitectura de avaliação, não num classificador novo: elegibilidade temporal, apuração do desfecho, divisão por paciente, operações ajustadas só no treino, selecção específica da representação, agregação de probabilidades e comparações com correcção de multiplicidade formam uma cadeia auditável.
	Esse desenho permite comparar, no mesmo estimando, \code{t1\_d21\_d32} versus \code{t1\_only} em cada família, a união tardia de cinco SVMs, e o conjunto de atributos que \code{l1\_stable} retém em cada dobra.
	
	\paragraph{Perguntas.}
	No hipocampo, para sMCI versus pMCI:
	(1)~Q4 (\code{t1\_d21\_d32}: T1, D21, D32) agrega AUC relativamente a \code{t1\_only} \emph{dentro} de cada família?
	(2)~quais atributos \code{l1\_stable} retém (frequência em $50$ avaliações externas), em vez de uma assinatura única?
	(3)~a fusão tardia (três especificações pré-especificadas) supera o teto unimodal?
	Quantifica-se $\Delta\mathrm{AUC}=\mathrm{AUC}(\mathrm{Q4})-\mathrm{AUC}(\code{t1\_only})$ nas cinco famílias e nas quatro células do fatorial janela clínica $\times$ intervalo entre imagens.
	
	Contribuições para avaliação reprodutível em neuroimagem longitudinal:
	\begin{itemize}
		\item protocolo que separa visitas preditoras em estágio MCI da apuração do desfecho e documenta a fronteira de ajuste de cada operação dependente da amostra sob validação aninhada no paciente;
		\item matriz controlada em cinco famílias hipocampais, \code{t1\_only} versus Q4, quatro coortes e fusão tardia (sem concatenação precoce da união), sob o mesmo estimando de AUC no paciente e o mesmo SVM;
		\item leitura por papéis: a forma satura no $t_0$; volume, textura e primeira ordem agregam no tempo; a união tardia é o contraste que passa o teto unimodal;
		\item evidência de apoio à decisão: variáveis clínicas baratas permanecem competitivas nesta coorte; a RM hipocampal é adjunto estrutural cuja contribuição incremental deve ser demonstrada, não presumida.
	\end{itemize}
	
	O desfecho primário é a classificação binária sMCI versus pMCI.
	O restante do artigo organiza-se assim: a Secção~\ref{sec:metodos} descreve coorte, extração de atributos e validação; a Secção~\ref{sec:resultados} reporta os resultados prognósticos e comparativos; a Secção~\ref{sec:discussao} discute achados e limitações; a Secção~\ref{sec:conclusoes} conclui.
	
	
\subsection{Trabalhos relacionados}
\label{subsec:relacionados}

A previsão automática de estabilidade ou conversão do MCI foi estudada com conjuntos de atributos, esquemas de validação e janelas de desfecho muito distintos.
Na investigação prognóstica baseada na ADNI, MCI estável e progressivo aparecem também como ncMCI/cMCI ou MCInc/MCIc, clinicamente análogos a sMCI e pMCI.
O desempenho depende, portanto, não só do classificador, mas de os preditores serem clínicos, estruturais, multimodais, regionais ou de cérebro inteiro, e de a informação temporal limitar-se ao baseline ou cobrir várias visitas de seguimento.

Vários estudos mostram que o acompanhamento cognitivo denso pode ser informativo mesmo sem RM.
\citet{bucholc-2023} propuseram um ensemble híbrido no National Alzheimer's Coordinating Center (NACC) para prever conversão MCI$\to$demência a partir de trajectórias cognitivas e funcionais.
A acurácia de classificação atingiu $85{,}0$--$87{,}5\%$, mas a AUC do ensemble no teste retido ficou perto de $0.60$, ao passo que florestas aleatórias chegaram a AUC até $0.76$ quando havia quatro visitas.
\citet{grassi-2019} previram conversão a 3 anos na ADNI só com variáveis sociodemográficas e neuropsicológicas e reportaram AUROC $0.88$ (IC bootstrap 95\%: $0.85$--$0.91$), com sensibilidade $77{,}7\%$ e especificidade $79{,}9\%$ no ponto de melhor acurácia balanceada.
Esses resultados indicam que o rastreio clínico carrega sinal prognóstico substancial e constitui um comparador forte para modelos só de imagem.

Estudos de RM estrutural recorrem frequentemente a representações de cérebro inteiro de alta dimensão ou a fusão com escores cognitivos.
\citet{moradi-2015} aplicaram \emph{low-density separation} a T1 de baseline da ADNI para previsão precoce de conversão em MCI, com AUC $0.766$ em morfometria voxel-a-voxel isolada e AUC $0.902$ ao combinar RM com um agregado cognitivo de baseline.
\citet{ye-2012} usaram selecção estável esparsa (Biosignature-15) para prever conversão MCI$\to$AD em quatro anos, com AUC $0.859$ em validação leave-one-out (LOOCV) na ADNI de baseline.
\citet{casanova-2013} combinaram atributos de RM em larga escala, baseados em tecido, com informação cognitiva num quadro de avaliação de risco por aprendizagem automática na ADNI.
Tais estudos mostram que AUC altas surgem sobretudo com desenhos de cérebro inteiro ou multimodais, e com definições de desfecho que podem diferir da discriminação sMCI/pMCI em horizonte curto e regras estritas.

Quando o âmbito espacial se reduz a uma única ROI, o desempenho em geral cai.
\citet{cuingnet-2011} compararam dez métodos de classificação na ADNI e reportaram que a discriminação conversor versus não conversor, inclusive com atributos de volume hipocampal, não excedeu o acaso de forma significativa.
\citet{guo-2017a} estudaram espessura cortical relacionada com conversão para sMCI versus cMCI, com acurácia $98{,}84\%$ e sensibilidade/especificidade $97{,}5/100$ em LOOCV, mas sem reportar AUC.
Esses achados são coerentes com a ideia de que restringir a análise a uma região focal --- hipocampo ou córtex --- produz desempenho prognóstico mais conservador do que modelos voxel de cérebro inteiro.

Trabalho mais recente incorpora RM temporal via aprendizagem de dicionários, redes profundas ou radiómica.
\citet{zhang-2019} analisaram RM estrutural longitudinal para classificação AD/MCI, \citet{lin-2020} usaram aprendizagem de dicionários longitudinal para prever progressão MCI$\to$AD, e \citet{aghajanian-2025} combinaram aprendizagem profunda com radiómica de RM estrutural longitudinal para previsão de progressão.
Essas abordagens mostram interesse crescente em modelação temporal, mas diferem em definição de coorte, controlo de vazamento e comparação com baselines clínicos.

Neste contexto, o presente trabalho não visa reproduzir o estado da arte de cérebro inteiro multimodal.
Visa especificar e estressar um pipeline prognóstico focado no hipocampo na ADNI, com configuração primária pré-especificada (SVM, \code{l1\_stable}, AUC no paciente) e ablações sistemáticas: cinco famílias de atributos, \code{t1\_only} versus Q4, quatro células janela$\times$intervalo, e fusão tardia em três especificações pré-especificadas.
A contribuição original é o protocolo de avaliação coordenado: todas as representações candidatas correm sob elegibilidade temporal explícita, validação aninhada no paciente, ajuste fold-wise e o mesmo estimando OOF.
Ablações negativas e positivas tornam-se evidência de desenho para futuros sistemas de apoio à decisão, em vez de tratar a maior AUC observada como único resultado.
Sob esse protocolo, a forma no $t_0$ atinge AUC paciente $0.761\pm0.043$ em \code{48m\_12m}; a textura Q4 empata esse teto ($0.762\pm0.044$); a fusão tardia Q4/âncora chega a $0.802\pm0.040$; os escores clínicos de $t_0$ permanecem competitivos ($0.812\pm0.040$).
A Tabela~\ref{tab:literature_comparison} resume literatura representativa ao lado destes resultados, sob definições de tarefa heterogéneas mas clinicamente análogas.
	
	% =============================================================================
	\section{Material e métodos}
	\label{sec:metodos}
	
	\subsection{Definição da coorte}
	\label{sec:populacao}
	
	
	Os dados utilizados neste estudo foram obtidos da base de dados da Alzheimer's Disease Neuroimaging Initiative (ADNI)~\cite{petersen-2009,jr-2008}.
	A ADNI foi lançada em 2003 como parceria público--privada para avaliar marcadores clínicos, genéticos, bioquímicos e de neuroimagem longitudinais da doença de Alzheimer (AD).
	Este estudo utilizou RM estrutural T1-ponderada adquirida com sequências tridimensionais Magnetization-Prepared Rapid Gradient-Echo (MP-RAGE) das fases ADNI-1, ADNI-2 e ADNI-GO, coletadas em centros nos Estados Unidos e no Canadá em scanners de 1.5~T e 3~T da GE Healthcare, Philips Medical Systems e Siemens Medical Solutions.
	
	A coleção ADNI inicial compreendia 15{,}016 imagens de 1{,}495 participantes.
	Duas imagens de um participante com designação de sexo inválida foram excluídas, restando 15{,}014 imagens de 1{,}494 participantes.
	
	O controle de qualidade foi aplicado a esse conjunto completo de volumes T1, ainda sem rótulos diagnósticos e antes da construção das coortes.
	Métricas de qualidade de imagem (IQMs) de referência livre foram calculadas com \code{MRQy}~\citep{sadri-2020}.
	Candidatos a outliers foram rastreados com as métricas \code{VRX}, \code{COL}, \code{NUM}, \code{CPP}, \code{PSNR}, \code{SNR7}, \code{CNR}, \code{CVP}, \code{CJV}, \code{EFC} e \code{FBER}, escolhidas pela relevância a recorte de FOV, geometria de aquisição, ruído e degradação relacionada a movimento ou artefato.
	Para cada métrica, outliers foram definidos pelas cercas de Tukey~\citep{tukey-1977},
	\begin{equation}
		\left[Q_1 - \alpha\,\mathrm{IQR},\, Q_3 + \alpha\,\mathrm{IQR}\right],
		\qquad \mathrm{IQR} = Q_3 - Q_1,
		\label{eq:tukey}
	\end{equation}
	com $\alpha = 1.5$ e direcionalidade específica da métrica (valores baixos indesejáveis, altos indesejáveis, ou bilateral).
	Cada volume recebeu uma pontuação de não conformidade igual ao número de IQMs sinalizadas; volumes com pontuação $\geq 3$ entraram em uma lista curta ($69/15{,}014$; $0.46\%$).
	Os quartis e os limites de IQR foram estimados uma única vez a partir desse conjunto de 15{,}014 volumes, sem rótulos diagnósticos, antes da construção das coortes e da validação cruzada; não foram recalculados nas partições externas.
	Esse rastreio ao nível do conjunto de dados foi, portanto, distinto da modelagem ajustada por partição.
	Todos os casos da lista curta foram inspecionados visualmente no ITK-SNAP~\citep{yushkevich-2006} (v4.2.0).
	A inspeção confirmou 12 identificadores de aquisição para exclusão por FOV incompleto, aquisições danificadas ou faixas de ruído estruturado.
	
	Em seguida procedeu-se à construção da coorte clínica.
	Na ADNI, o diagnóstico associado à primeira visita (baseline) não acompanha necessariamente a progressão longitudinal; por isso, cada exame de RM foi alinhado ao diagnóstico clínico e às escalas cognitivas/funcionais (ADAS, CDR, MMSE, FAQ) por correspondência temporal, exigindo distância absoluta $\leq 3$ meses ($\pm 90$ dias) entre a data de aquisição e a data da avaliação clínica, de modo que o rótulo corresponda o mais possível ao momento da imagem.
	Assim, imagens do mesmo paciente podem receber rótulos clínicos distintos ao longo do tempo.
	Antes do pareamento imagem--diagnóstico, os registros diagnósticos longitudinais foram filtrados: removeram-se demências diferentes da doença de Alzheimer, diagnósticos ou datas em falta, e reversões clinicamente implausíveis nas fases contínuas da AD --- MCI$\rightarrow$CN ou AD$\rightarrow$MCI/CN.
	Esse rastreio identificou 40 participantes com reversão a partir de AD e, após a sua remoção, mais 128 com reversão MCI$\rightarrow$CN, totalizando 168 participantes únicos excluídos; destes, 112 estavam representados na coleção de imagem e correspondiam a 1{,}325 imagens.
	O pareamento RM--diagnóstico na janela de $\pm 90$ dias excluiu mais 666 imagens, e a exigência de escores clínicos completos excluiu outras 92, resultando em 12{,}931 imagens rotuladas de 1{,}304 participantes.
	Dos 12 identificadores rejeitados pelo controle visual de qualidade, 10 ainda estavam presentes após o pareamento clínico e foram removidos nesta etapa; os outros dois já tinham sido eliminados pelos filtros diagnósticos precedentes.
	A tabela após o controle de qualidade continha, portanto, 12{,}921 imagens.
	A remoção de 3{,}757 aquisições cuja descrição ADNI as identificava como repetições (\emph{repeats}) reduziu o conjunto a 9{,}164 candidatas longitudinais utilizáveis.
	Esse pool de 9{,}164 imagens é o ponto de partida para a formação dos conjuntos longitudinais descritos a seguir; ainda não coincide com a coorte final de análise (Tabela~\ref{tab:adni_study}), que resulta da exigência de exatamente três imagens por paciente sob as regras de janela clínica e intervalo entre aquisições.
	
	
	\subsubsection{Formação dos conjuntos}
	\label{subsec:formacao}
	
	A formação de cada conjunto longitudinal segue duas etapas em sequência, para que o rótulo clínico não dependa de qual imagem acabou sendo escolhida como preditora.
	
	\paragraph{Etapa 1 --- rótulo clínico.}
	Para cada paciente, ordenam-se todas as suas RM pela data de aquisição.
	A data da primeira RM define o instante de referência $t_0$ (baseline).
	Sobre esse instante fixa-se uma \textbf{janela de observação clínica} de 36 ou 48 meses, com término em $t_{\mathrm{end}} = t_0 +$\,janela.
	Olhando toda a trajetória diagnóstica do paciente (não apenas três imagens), decide-se se ele é CN, sMCI, pMCI, AD, ou se deve ser excluído.
	Só depois dessa decisão passa-se à escolha das imagens que alimentarão o modelo.
	
	\paragraph{Etapa 2 --- três imagens preditoras.}
	Cada paciente contribuirá com exatamente três RM preditoras
	\begin{equation}
		i=\{i_1,\,i_2,\,i_3\},
		\quad t(i_1)<t(i_2)<t(i_3),
	\end{equation}
	isto é, uma imagem inicial, uma intermediária e uma final.
	Três é o menor número que ainda permite descrever uma trajetória temporal (início, meio e fim) sem exigir séries longas, raras em repositórios públicos.
	A escolha é cronológica e usa um intervalo nominal entre aquisições (6 ou 12 meses) com tolerância $\pm 2$ meses:
	\begin{enumerate}
		\item $i_1$ é a primeira RM elegível no baseline ($t_0$);
		\item $i_2$ é a \textbf{próxima} RM, em ordem de data, cujo intervalo em relação a $i_1$ cai dentro da faixa de tolerância;
		\item $i_3$ é a \textbf{próxima} RM após $i_2$ cujo intervalo em relação a $i_2$ cai na \textbf{mesma} faixa.
	\end{enumerate}
	Não se trata de ``três imagens \emph{depois} do baseline'' (o que seriam quatro visitas): são três no total --- o baseline e duas visitas seguintes que respeitam o espaçamento.
	
	As faixas adotadas ($\pm 2$ meses) são disjuntas e evitam ambiguidade no bordo dos 9 meses.
	Com tolerância $\pm 3$, o intervalo de 6 meses cobriria $[3,9]$ e o de 12 meses cobriria $[9,15]$, compartilhando o mês 9: uma imagem adquirida nove meses após a anterior poderia pertencer a qualquer um dos dois esquemas.
	Com $\pm 2$:
	
	\begin{table}[htbp]
		\centering
		\caption{Intervalos nominais entre aquisições e faixas de tolerância.}
		\label{tab:bandas}
		\begin{tabular}{ll}
			\toprule
			Intervalo nominal & Faixa admitida (meses) \\
			\midrule
			6\,m  & $[4,\ 8]$ \\
			12\,m & $[10,\ 14]$ \\
			\bottomrule
		\end{tabular}
	\end{table}
	
	Combinando as duas durações de janela clínica com os dois intervalos entre imagens obtêm-se \textbf{quatro coortes} de análise: 36\,meses com intervalo de 6\,meses, 36\,meses com intervalo de 12\,meses, 48\,meses com intervalo de 6\,meses e 48\,meses com intervalo de 12\,meses (notação \code{36m\_6m}, \code{36m\_12m}, \code{48m\_6m}, \code{48m\_12m}).
	
	\paragraph{Definição dos grupos clínicos.}
	\begin{itemize}
		\item \textbf{CN:} toda a trajetória diagnóstica é exclusivamente CN, e existe pelo menos uma visita com data $\ge t_{\mathrm{end}}$ que confirma a estabilidade após o fim da janela.
		As três preditoras $i_1,i_2,i_3$ são então escolhidas entre imagens CN até $t_{\mathrm{end}}$, respeitando a faixa de intervalo.
		
		\item \textbf{sMCI:} o baseline é MCI, não há diagnóstico de AD até $t_{\mathrm{end}}$, existe pelo menos uma visita com data $\ge t_{\mathrm{end}}$, e a \textbf{primeira} dessas visitas pós-janela permanece MCI.
		Se a primeira visita após $t_{\mathrm{end}}$ for AD (ou outro diagnóstico não-MCI), o paciente é excluído: o desfecho ficou ``no meio'' da margem de confirmação e não se classifica de forma limpa.
		Se não houver visita após $t_{\mathrm{end}}$, também é excluído por falta de confirmação.
		Essa definição \textbf{não} afirma que o paciente nunca converterá a AD; exige apenas ausência de conversão no horizonte pré-especificado, com confirmação na primeira avaliação após o fim da janela.
		A visita de confirmação \textbf{não} precisa respeitar a faixa $\pm 2$ meses --- essa faixa serve apenas para escolher $i_1,i_2,i_3$ (todas MCI, até $t_{\mathrm{end}}$).
		
		\item \textbf{pMCI:} o baseline é MCI e ocorre pelo menos um diagnóstico de AD com data $\le t_{\mathrm{end}}$.
		Preferencialmente selecionam-se três preditoras ainda em estágio MCI, antes do primeiro AD ($\mathrm{MCI},\mathrm{MCI},\mathrm{MCI}$), desde que exista conversão documentada dentro da janela.
		Quando não há três aquisições MCI com o espaçamento exigido antes da conversão, admite-se uma regra de inclusão ampliada: $\mathrm{MCI},\mathrm{MCI},\mathrm{AD}$, em que $i_3$ é o \textbf{primeiro} AD (único AD no conjunto preditor).
		Não se admite $\mathrm{MCI},\mathrm{AD},\mathrm{AD}$: apenas o primeiro AD define o desfecho; um segundo AD no preditor afastaria o último MCI do horizonte e misturaria alterações estruturais já pós-conversão.
		
		\item \textbf{AD:} toda a trajetória é exclusivamente AD; as três preditoras são AD, escolhidas na faixa de intervalo até $t_{\mathrm{end}}$.
	\end{itemize}
	
	Em resumo: primeiro classifica-se o paciente pela história clínica na janela; depois escolhem-se, em ordem temporal, três imagens cujo espaçamento cai na faixa de tolerância.
	Pacientes que não fornecem três imagens elegíveis sob essas regras ficam de fora da coorte correspondente, mesmo que o rótulo clínico esteja definido.
	
	
	\subsubsection{Coortes de estudo}
	\label{subsec:coortes}
	
	A Tabela~\ref{tab:coortes} resume as contagens sMCI/pMCI após a formação dos conjuntos em cada combinação janela$\times$intervalo.
	
	\begin{table}[htbp]
		\centering
		\caption{Contagens sMCI/pMCI em conjuntos (tolerância $\pm 2$ meses).}
		\label{tab:coortes}
		\begin{tabular}{lll}
			\toprule
			Janela & Intervalo entre imagens & sMCI / pMCI \\
			\midrule
			36\,m & $6\pm 2\ [4,8]$   & 125 / 106 \\
			36\,m & $12\pm 2\ [10,14]$ & 121 / 33 \\
			48\,m & $6\pm 2\ [4,8]$   & 73 / 120 \\
			48\,m & $12\pm 2\ [10,14]$ & \textbf{72 / 48} \\
			\bottomrule
		\end{tabular}
	\end{table}
	
	As 9{,}164 candidatas longitudinais não coincidem com o total da Tabela~\ref{tab:adni_study}: esta última é a \textbf{união} das imagens únicas presentes em pelo menos uma das quatro coortes de análise (CN, sMCI, pMCI e AD), após a exigência de exatamente três aquisições e sem duplicar a mesma imagem quando ela entra em mais de uma combinação janela$\times$intervalo --- 2{,}038 imagens de 560 participantes, estratificadas por fase ADNI de aquisição.
	
	A Tabela~\ref{tab:clin_cohorts} resume demografia e escores cognitivos/funcionais nas três visitas preditoras ($t_0$, $t_1$, $t_2$) para sMCI e pMCI nas quatro coortes.
	CN/AD e CDR global foram omitidos (continuum de sanidade / constante 0{,}50 no baseline MCI).
	Valores contínuos: média $\pm$ desvio-padrão; $n$ apenas na coluna de pacientes.
	
	\begin{table}[htbp]
		\centering
		\footnotesize
		\caption{Características clínicas sMCI/pMCI nas quatro coortes (tolerância $\pm 2$; inclusão ampliada de pMCI ativa).}
		\label{tab:clin_cohorts}
		\resizebox{\textwidth}{!}{%
			\begin{tabular}{lllrrccccc}
				\toprule
				Coorte & Grupo & Slot & $n$ & Sexo M/F & Idade & MMSE & ADAS & CDR-SB & FAQ \\
				\midrule
				\multirow{6}{*}{\code{36m\_6m}}
				& \multirow{3}{*}{sMCI}
				& $t_0$ & 125 & 80/45 & $72.63\pm7.18$ & $27.78\pm1.75$ & $9.05\pm3.61$ & $1.22\pm0.64$ & $1.93\pm3.09$ \\
				& & $t_1$ & 125 & 80/45 & $73.16\pm7.16$ & $28.04\pm1.77$ & $8.79\pm3.64$ & $1.23\pm0.77$ & $2.22\pm3.33$ \\
				& & $t_2$ & 125 & 80/45 & $73.67\pm7.19$ & $27.99\pm1.88$ & $8.47\pm4.11$ & $1.36\pm0.95$ & $2.53\pm3.56$ \\
				\cmidrule{2-10}
				& \multirow{3}{*}{pMCI}
				& $t_0$ & 106 & 56/50 & $74.39\pm7.33$ & $26.64\pm1.74$ & $13.54\pm4.49$ & $1.94\pm1.00$ & $5.34\pm4.42$ \\
				& & $t_1$ & 106 & 56/50 & $74.90\pm7.26$ & $25.32\pm2.65$ & $14.84\pm5.00$ & $2.44\pm1.01$ & $6.79\pm4.75$ \\
				& & $t_2$ & 106 & 56/50 & $75.44\pm7.30$ & $24.68\pm2.70$ & $15.75\pm5.83$ & $3.19\pm1.41$ & $8.95\pm5.66$ \\
				\midrule
				\multirow{6}{*}{\code{36m\_12m}}
				& \multirow{3}{*}{sMCI}
				& $t_0$ & 121 & 80/41 & $73.15\pm7.06$ & $27.80\pm1.74$ & $9.05\pm3.60$ & $1.21\pm0.63$ & $1.85\pm3.01$ \\
				& & $t_1$ & 121 & 80/41 & $74.21\pm7.07$ & $28.02\pm1.87$ & $8.62\pm4.20$ & $1.33\pm0.88$ & $2.58\pm3.48$ \\
				& & $t_2$ & 121 & 80/41 & $75.22\pm7.05$ & $27.85\pm2.17$ & $9.08\pm4.54$ & $1.51\pm1.20$ & $2.70\pm3.65$ \\
				\cmidrule{2-10}
				& \multirow{3}{*}{pMCI}
				& $t_0$ & 33 & 22/11 & $74.30\pm7.35$ & $27.00\pm1.62$ & $12.30\pm2.85$ & $2.09\pm1.17$ & $5.45\pm4.90$ \\
				& & $t_1$ & 33 & 22/11 & $75.33\pm7.34$ & $25.88\pm2.13$ & $13.31\pm4.07$ & $2.67\pm1.19$ & $7.55\pm4.94$ \\
				& & $t_2$ & 33 & 22/11 & $76.39\pm7.34$ & $23.33\pm2.91$ & $17.24\pm4.78$ & $4.45\pm1.40$ & $12.76\pm6.02$ \\
				\midrule
				\multirow{6}{*}{\code{48m\_6m}}
				& \multirow{3}{*}{sMCI}
				& $t_0$ & 73 & 48/25 & $72.29\pm6.80$ & $27.89\pm1.65$ & $8.80\pm3.53$ & $1.11\pm0.61$ & $1.88\pm3.32$ \\
				& & $t_1$ & 73 & 48/25 & $72.84\pm6.74$ & $28.30\pm1.40$ & $8.33\pm3.40$ & $1.09\pm0.71$ & $2.07\pm3.40$ \\
				& & $t_2$ & 73 & 48/25 & $73.30\pm6.81$ & $28.16\pm1.84$ & $7.83\pm3.49$ & $1.16\pm0.95$ & $2.30\pm3.36$ \\
				\cmidrule{2-10}
				& \multirow{3}{*}{pMCI}
				& $t_0$ & 120 & 62/58 & $74.35\pm7.20$ & $26.65\pm1.71$ & $13.20\pm4.41$ & $1.89\pm0.98$ & $5.17\pm4.46$ \\
				& & $t_1$ & 120 & 62/58 & $74.85\pm7.15$ & $25.43\pm2.60$ & $14.44\pm4.99$ & $2.37\pm1.00$ & $6.46\pm4.81$ \\
				& & $t_2$ & 120 & 62/58 & $75.40\pm7.18$ & $24.89\pm2.65$ & $15.38\pm5.65$ & $3.04\pm1.43$ & $8.53\pm5.81$ \\
				\midrule
				\multirow{6}{*}{\textbf{\code{48m\_12m}}}
				& \multirow{3}{*}{sMCI}
				& $t_0$ & 72 & 49/23 & $72.88\pm6.70$ & $28.00\pm1.63$ & $8.71\pm3.40$ & $1.07\pm0.59$ & $1.85\pm3.30$ \\
				& & $t_1$ & 72 & 49/23 & $73.93\pm6.73$ & $28.32\pm1.69$ & $8.07\pm3.65$ & $1.17\pm0.89$ & $2.43\pm3.37$ \\
				& & $t_2$ & 72 & 49/23 & $74.96\pm6.68$ & $28.29\pm1.86$ & $8.65\pm3.59$ & $1.29\pm1.03$ & $2.26\pm3.26$ \\
				\cmidrule{2-10}
				& \multirow{3}{*}{pMCI}
				& $t_0$ & 48 & 28/20 & $73.98\pm6.83$ & $26.94\pm1.56$ & $12.18\pm2.95$ & $1.93\pm1.10$ & $4.83\pm4.88$ \\
				& & $t_1$ & 48 & 28/20 & $75.00\pm6.83$ & $25.90\pm1.95$ & $13.24\pm3.79$ & $2.45\pm1.18$ & $6.77\pm5.36$ \\
				& & $t_2$ & 48 & 28/20 & $76.06\pm6.86$ & $23.90\pm2.81$ & $16.44\pm4.81$ & $3.76\pm1.61$ & $10.56\pm6.51$ \\
				\bottomrule
			\end{tabular}%
		}
		\begin{flushleft}
			\footnotesize
			\textit{Nota:} Inclusão ampliada de pMCI (último preditor = primeiro AD): 46 em \code{36m\_6m}/\code{48m\_6m}; 33 em \code{36m\_12m}/\code{48m\_12m}.
			MMSE: Mini-Mental State Examination; ADAS: Alzheimer's Disease Assessment Scale; CDR-SB: Clinical Dementia Rating -- Sum of Boxes; FAQ: Functional Activities Questionnaire.
			As quatro coortes entram como células do fatorial janela$\times$intervalo; \code{36m\_6m} tem o maior $n$.
		\end{flushleft}
	\end{table}
	
	\begin{table}[htbp]
		\centering
		\footnotesize
		\caption{Imagens T1w MP-RAGE por fase ADNI (união das quatro coortes de análise; imagens únicas por identificador de aquisição).}
		\label{tab:adni_study}
		\resizebox{\textwidth}{!}{%
			\begin{tabular}{lrrcccc}
				\toprule
				Fase ADNI & Imagens & Pacientes & Sexo (M/F) & Idade (anos) & Campo (1.5T / 3T) & Fabricante (n) \\
				\midrule
				ADNI~1  & 1349 & 370 & 203/167 & $76.14\pm6.91$ & 1261 / 88  & GE (768); Siemens (579); Philips (2) \\
				ADNI~GO & 51   & 27  & 16/11   & $70.61\pm6.76$ & 0 / 51     & Siemens (44); Philips (7) \\
				ADNI~2  & 638  & 190 & 101/89  & $73.63\pm7.02$ & 0 / 638    & Siemens (483); Philips (155) \\
				\midrule
				Total   & 2038 & 560 & 304/256 & $75.21\pm7.07$ & 1261 / 777 & GE (768); Siemens (1106); Philips (164) \\
				\bottomrule
			\end{tabular}%
		}
		\begin{flushleft}
			\footnotesize
			\textit{Nota:} Contagens sem duplicar a mesma imagem entre coortes (sobreposição janela$\times$intervalo).
			Idade = média $\pm$ DP no momento da aquisição.
			Sexo por paciente único dentro de cada fase (um paciente pode contribuir imagens a mais de uma fase; $370+27+190=587>560$).
			Modelos de equipamento (n imagens):
			ADNI~1 --- SIGNA EXCITE (445), Symphony (233), SIGNA HDx (196), Sonata (146), GENESIS\_SIGNA (123), Avanto (100), Trio (34), SymphonyTim (25), Allegra (17), SonataVision (13), TrioTim (11);
			ADNI~GO --- Verio (25), TrioTim (18), Intera (4), Achieva (3), Skyra (1);
			ADNI~2 --- TrioTim (297), Achieva (95), Verio (95), Skyra (90), Intera (38), Ingenia (17).
		\end{flushleft}
	\end{table}
	
	Também se exploraram as bases \textbf{AIBL} e \textbf{OASIS} sob a mesma lógica de elegibilidade, mas sem número suficiente de conjuntos longitudinais com três imagens.
	
	A tarefa principal é a classificação \textbf{sMCI versus pMCI}.
	O mesmo pipeline (três T1 hipocampais, SVM, \code{l1\_stable}, AUC no paciente) aplica-se a \textbf{quatro cenários} do fatorial janela clínica $\times$ intervalo entre imagens --- não a uma coorte semente e três expansões.
	A janela (36 ou 48 meses) define o horizonte do rótulo; o intervalo nominal (6 ou 12 meses, $\pm 2$) define a cadência das preditoras.
	As quatro células entram no contraste unimodal Q4 versus \code{t1\_only} no mesmo nível: o objectivo é ver se o padrão de famílias se replica quando mudam horizonte e espaçamento, não coroar a célula com maior $\Delta$AUC.
	As populações sobrepõem-se e \textbf{não} são réplicas independentes; diferem $n$, balanço sMCI/pMCI e, nos intervalos de 12 meses, a fracção de pMCI com $i_3$ já em AD.
	Relato complementar (não substitui o fatorial): demografia sobretudo em \code{36m\_6m} (maior $n$); nomes de atributos, clínico, leaky e CN$\times$AD em \code{48m\_12m} (palco das ablações que não se repetem nas quatro células).
	A fusão tardia reporta-se nas \textbf{quatro} células.

	\paragraph{Janela de 36 versus 48 meses no mesmo intervalo entre imagens.}
	Não se deve esperar AUC idêntica em \code{36m\_6m} versus \code{48m\_6m} (nem em \code{36m\_12m} versus \code{48m\_12m}) apenas porque a representação \code{t1\_only} usa a visita de baseline.
	Na \textbf{interseção} de pacientes presentes nas duas coortes com o mesmo espaçamento nominal, o identificador de aquisição de $t_0$ coincide (e o triplo $t_0$--$t_1$--$t_2$ também), com o mesmo rótulo de grupo.
	Contudo as \textbf{populações não são as mesmas}: a duração da janela clínica altera a elegibilidade (quem entra ou sai).
	Exemplos no protocolo atual (tolerância $\pm 2$): sMCI 125 versus 73 (intervalo 6\,meses) e 121 versus 72 (intervalo 12\,meses), com dezenas de identificadores exclusivos de uma das janelas.
	Logo diferem $n$, balanço sMCI/pMCI, partições de validação cruzada e AUC --- inclusive em \code{t1\_only} --- mesmo com sobreposição de baselines.
	
	\subsection{Representação longitudinal}
	\label{subsec:q4}
	
	Cada sujeito contribui três aquisições T1 nas visitas $T_1$, $T_2$ e $T_3$ (intervalos nominais 6 ou 12 meses).
	O baseline \code{t1\_only} usa só os atributos de $T_1$.
	A representação longitudinal oficial \code{t1\_d21\_d32} (Q4) concatena \emph{níveis} das três visitas: colunas T1, D21 e D32, em que D21 e D32 são os mesmos descritores medidos em $T_2$ e $T_3$, não diferenças $T_b-T_a$.
	Não se usam \code{wide}/\code{abs}, \code{t1\_deltas} nem médias móveis como claim temporal: Q4 testa se \emph{estados} sucessivos acrescentam discriminação ao $t_0$, sem reparametrizar mudança.
	
	% =============================================================================
	\subsection{Pré-processamento}
	\label{sec:preprocessamento}
	
	Após obtenção dos dados brutos de RM na base ADNI, os volumes foram submetidos a um pipeline de pré-processamento próprio. Essas operações foram aplicadas de forma independente a cada imagem, com configurações de software fixas: modelos pré-treinados foram usados para extração de crânio e parcelamento; o denoising e a correção de campo de viés estimaram quantidades específicas de cada imagem; a normalização de intensidade e o alinhamento espacial usaram imagens de referência e parâmetros algorítmicos fixos. Nenhum parâmetro de pré-processamento foi estimado por agrupamento das imagens sMCI e pMCI.
	
%	Quando o NIfTI de entrada era tetradimensional, reteve-se um único volume 3D selecionando o primeiro quadro ao longo do quarto eixo com \code{ants.slice\_image(axis=3, idx=0)}. Essa conversão afetou 4 de 15{,}014 aquisições brutas (0{,}03\%); as quatro imagens afetadas continham dois quadros no quarto eixo. Nesta coorte T1 do ADNI, o quarto eixo correspondia a quadros repetidos ou extras, e não a uma pilha espectroscópica multi-eco; não se utilizou regra alternativa de seleção de quadro.
	
	Em seguida, todos os volumes ponderados em T1 passaram por uma cadeia fixa de pré-processamento no ecossistema ANTsX \cite{tustison-2021}, composta pelos seguintes passos sequenciais:
	
	\textbf{Extração de crânio:} Executada com um modelo de extração encefálica por aprendizagem profunda treinado em T1, da biblioteca ANTsPyNet (\code{antspynet.brain\_extraction}, \code{modality="t1"}) \cite{tustison-2021}.
		
	\textbf{Denoising:} Médias não locais adaptativas para ruído ricianiano \cite{manjon-2010}, implementadas com \code{ants.denoise\_image} usando \code{noise\_model="Rician"}, raio de patch $p=2$, raio de busca $r=3$ e \code{mask=None}. O parâmetro específico do ANTs \code{shrink\_factor=1} desativa o downsampling interno e é um controlo de implementação, não parte da formulação original.
		
	\textbf{Correção de campo de viés:} Correção N4ITK \cite{tustison-2010} com \code{ants.n4\_bias\_field\_correction}, usando \code{shrink\_factor=3}, \code{convergence=\{iters: [50, 50, 50, 50], tol:} \(1\times10^{-7}\)\code{\}} (quatro níveis de multi-resolução), \code{spline\_param=200}, \code{rescale\_intensities=True} e \code{mask=None}.
		
	\textbf{Alinhamento espacial:} As imagens clínicas foram registradas rigidamente ao espaço de referência MNI152 2009c \cite{fonov-2009,fonov-2011} com \code{ants.registration} do ANTsPy \cite{avants-2011a,tustison-2021}, corrigindo o alinhamento global e padronizando a orientação espacial.
	
	As imagens T1 utilizadas neste estudo já tinham sido submetidas, fora deste repositório, a extração de crânio, denoising por médias não locais e correção de campo de viés N4.
	O alinhamento espacial aqui aplicado consiste em registro rígido ao template MNI152 2009c~\cite{fonov-2011} em grelha isotrópica de $1{,}0$\,mm (\code{2\_resample.py}), com interpolação linear da intensidade e vizinho mais próximo para as máscaras de ROI, a segmentação tecidual e a máscara encefálica.
	
%	Não se aplicou \emph{histogram matching} neste fluxo.
	
	Os hipocampos esquerdo e direito (rótulos DKT 17 e 53) foram parcelados com \code{antspynet.desikan\_killiany\_tourville\_labeling} (\code{modality="t1"}, \code{do\_preprocessing=True}). A segmentação de classes teciduais --- líquido cefalorraquidiano (LCR), substância cinzenta (GM) e substância branca (WM) (rótulos 1, 2 e 3) --- foi obtida com \code{antspynet.deep\_atropos}, com a mesma modalidade e a mesma flag de pré-processamento.
	Embora a cadeia externa ANTsX já incluísse extração de crânio, denoising, N4, matching de histograma e alinhamento rígido ao MNI, reteve-se o pré-processamento interno da biblioteca para coincidir com o domínio esperado pelos pesos pré-treinados do ANTsPyNet. Os rótulos foram usados apenas para definição de ROI e fracções teciduais; as transformações internas do ANTsPyNet não foram tratadas como estágio treinável adicional.
	
	\subsection{Construção de templates groupwise estratificados}
\label{sec:groupwise}
	Para reduzir o viés de um único sujeito de referência \cite{reuter-2012} e fornecer um espaço normativo demografamente adequado, construíram-se templates volumétricos médios a partir de imagens T1 pré-processadas de indivíduos cognitivamente normais (CN). Os templates foram estratificados por sexo e década de idade (50{,}9--59{,}9, 60{,}9--69{,}9, \ldots, 90{,}9--99{,}9 anos) e serviram como referências de envelhecimento saudável para o registro deformável visita a visita.
	Para cada estrato sexo--década, reteve-se até $N_{\max}=20$ aquisições CN. Quando havia mais de 20 exames CN elegíveis, aplicou-se amostragem aleatória sem reposição com semente fixa \code{RANDOM\_SEED=7}, para reprodutibilidade. A escolha $N_{\max}=20$ equilibra representatividade demográfica e custo computacional.
	A amostra de referência CN foi selecionada uma única vez, antes da validação cruzada, e a seleção não foi repetida nos folds externos. Nenhum participante sMCI ou pMCI contribuiu para a construção dos templates. Para o desfecho primário sMCI/pMCI, os templates resultantes foram portanto tratados como referências normativas fixas, provenientes de um grupo disjunto de ambas as classes prognósticas, e não como transformações ajustadas por fold.
	O alinhamento groupwise começou com um pré-registro global mascarado de cada imagem selecionada a um alvo comum, aplicando transformações rígidas seguidas de afins. O template inicial $T^{(0)}$ foi definido como a média aritmética dessas imagens pré-alinhadas, reamostradas no espaço do alvo. Seguiu-se refinamento não linear por Symmetric Normalization (SyN) \cite{avants-2014} ao longo de cinco iterações: em cada iteração, todas as imagens do estrato foram registradas de forma não linear ao template corrente sob máscaras encefálicas, e o template foi actualizado pela média das imagens deformadas. Após a última iteração SyN, o template groupwise foi pós-processado para consistência de atlas: recomputou-se uma máscara encefálica sobre o template médio, e as intensidades intracranianas foram reescaladas de forma robusta, recortando aos percentis 2--98 e mapeando linearmente esse intervalo para $[0,1]$. Este passo padroniza o contraste visual entre estratos e não altera as transformações deformáveis usadas no pipeline prognóstico.

	\subsection{Registro normativo visita a visita}\label{subsec:deslocamento_longitudinal}

	Os atributos de deslocamento derivam de \emph{registro normativo visita a visita}, não de mapas de diferença inter-visita. Para cada participante, o template groupwise CN foi fixado pelo sexo e pela década de idade na primeira visita de entrada do modelo ($T_1$), independentemente do diagnóstico clínico corrente \cite{wyman-2013}. Essa atribuição de template ancorada em $T_1$ manteve-se constante em $T_2$ e $T_3$, independentemente da progressão clínica subsequente.
	Em cada visita $t \in \{T_1,T_2,T_3\}$, realizou-se um registro SyN independente no espaço de $1{,}0$\,mm, com o volume T1 clínico pré-processado como imagem \textbf{fixa} e o template CN atribuído como imagem \textbf{móvel} (\code{ants.registration}). O ANTs estimou, portanto, uma transformação template$\to$clínico $F_t$ (móvel $\rightarrow$ fixa) e a sua inversa $I_t$ (fixa $\rightarrow$ móvel). Affine, warp directo e warp inverso foram armazenados por imagem.
	A informação longitudinal entra no classificador pela representação Q4: mapas escalares dos três registros template$\to$clínico separados concatenam-se como slots T1, D21 e D32.
	\textbf{Não} se compuseram warps entre visitas (p.ex.\ $I_{T_2}\circ F_{T_1}$) e \textbf{não} se calcularam campos de deslocamento baseline$\to$follow-up para os atributos prognósticos reportados.
	Cada visita quantifica o desvio anatómico relativamente a uma referência normativa CN comum naquele instante, não uma mudança morfológica inter-visita explícita.
	

	\subsection{Extração de atributos}\label{sec:feature-extraction}

	Por aquisição extraem-se cinco famílias: volume, forma, textura GLCM, primeira ordem e deslocamento normativo.
	O armazém guarda o bloco IBSI completo e os mapas do DVF.
	O classificador recebe a \emph{classe} de cada família após uma denylist só de definição (\code{keep\_*\_feat} em \code{ablation\_prep.py}; Tabela~\ref{tab:dims}): não há lista permitida de 4--5 nomes.
	A selecção amostral (variância, correlação, pool \code{l1\_stable}) corre depois, só no treino externo.

	\subsubsection{Atributos volumétricos}

	Por ROI intersectada com a máscara encefálica, \code{3\_feat\_vol.py} regista volumes em mm$^3$ e fracções \emph{dentro da ROI}: $\mathtt{gm\_norm}=V_{\mathrm{GM}}/V_{\mathrm{ROI}}$ (análogo para WM e LCR).
	O ICV (\code{ICV\_mask\_mm3}) é o volume da máscara encefálica (linha \code{\_\_global\_\_}).
	No pós-processamento, descritores de tamanho (mm$^3$ e malha/eixos do PyRadiomics) dividem-se por esse ICV; as fracções \textbf{não} se renormalizam pelo ICV --- não são volumetria clássica em mm$^3$ \cite{wyman-2013}.
	Entram \code{gm\_norm}, \code{wm\_norm}, \code{csf\_norm} e \code{original\_shape\_MeshVolume} (já /ICV).
	Saem volumes crus em mm$^3$.
	O volume de malha fica no bloco de volume, não no de forma.

	\subsubsection{Atributos radiómicos (PyRadiomics)}\label{sec:radiomics}

	Forma, textura e primeira ordem extraem-se com PyRadiomics \code{3.1.1.dev111+g8ed579383} \cite{griethuysen-2017}.
	Nomes e definições seguem a terminologia IBSI da biblioteca \cite{zwanenburg-2020}; não se reivindica validação IBSI independente.

	A extração é independente por ROI e aquisição (\code{3\_feat\_rad.py}), sobre volumes já rígidos ao MNI152 2009c em grelha isotrópica de $1{,}0$\,mm.
	O PyRadiomics não reamostra (\code{resampledPixelSpacing} omitido).
	Grelha e afinidade imagem/máscara verificam-se com tolerância absoluta \(10^{-5}\).
	Cada rótulo intersecta a máscara encefálica; ROI ausente grava-se NaN.
	Só o tipo \code{Original} (sem LoG nem wavelet).
	Discretização: \code{binCount}$=64$, para evitar o default CT \code{binWidth}$=25$.

	Por ROI e visita o armazém tem 107 descritores (18 de primeira ordem, 14 de forma, 24 GLCM, 16 GLRLM, 16 GLSZM, 5 NGTDM, 14 GLDM).
	Nas ablações isoladas do hipocampo entram, após denylist:
	\begin{itemize}
		\item forma: os 12 \code{original\_shape\_*} excepto \code{MeshVolume} e \code{VoxelVolume} (tamanho já está no bloco volume);
		\item textura: os 24 \code{original\_glcm\_*} (GLCM Original); saem GLRLM/GLSZM/GLDM/NGTDM e filtros wavelet/LoG;
		\item primeira ordem: os 16 \code{original\_firstorder\_*} excepto \code{Energy} e \code{TotalEnergy} (escala com o número de voxels).
	\end{itemize}
	Contagens L$+$R \emph{antes} do L1: Tabela~\ref{tab:dims}.

	\subsubsection{Atributos de deslocamento}\label{sec:displacement-features}

	A referência é o atlas groupwise CN do sexo e da década em $T_1$, independente do diagnóstico \cite{wyman-2013}, e mantém-se em $T_2$ e $T_3$.
	Em cada visita, SyN com template CN \textbf{fixo} e T1 clínico \textbf{móvel} (\code{3\_feat\_gen\_dvf\_v2.py}).
	Os mapas vivem no domínio do template. O campo obtém-se aplicando $I_t=[\text{affine},\text{inverse warp}]$ às coordenadas físicas do template (fixed $\rightarrow$ moving):
	\begin{equation}
		\mathbf{u}_t(x)=I_t(x)-x,\qquad x\in\Omega_{\mathrm{CN}}.
		\label{eq:dvf-template}
	\end{equation}
	ROIs e máscara encefálica do sujeito transportam-se ao template pela cadeia directa (Warp$+$Affine, vizinho mais próximo).
	De $\mathbf{u}_t$ derivam-se $\log J$, $\|\mathbf{u}\|$, divergência, $(u_x,u_y,u_z)$, rotacional, escalares de Green--Lagrange e $\|\varepsilon\|_F$ da deformação infinitesimal $\varepsilon=\tfrac12(H+H^{\top})$ \cite{bouke-2023}.
	Não há mapas $T_a\rightarrow T_b$: $T_2$ e $T_3$ são desvio ao mesmo âncora CN, não deformação inter-visita.
	Por ROI hipocampal no template agregam-se média, desvio-padrão, percentis e, para $\log J$, $\|\mathbf{u}\|$ e $\|\varepsilon\|_F$, variância, assimetria e curtose.

	Entram momentos de \code{mag\_}, \code{logjac\_} e \code{strain\_fro\_} (15 sufixos por lado na representação longa).
	Saem componentes \code{ux}/\code{uy}/\code{uz}, sufixos \code{\_n} e percentis.
	\code{strain\_fro\_mean} resume a norma de Frobenius de Green--Lagrange; \code{strain\_fro\_variance} a da $\varepsilon$ infinitesimal --- tensores distintos.

	\subsection{Representações temporais}
	\label{sec:atributos}

	\begin{table}[htbp]
		\centering
		\caption{Representações temporais oficiais.}
		\label{tab:reps}
		\begin{tabular}{llp{7.4cm}}
			\toprule
			Nome & Papel & Definição \\
			\midrule
			\code{t1\_only} & baseline & só $t_0$ (coluna T1) \\
			\code{t1\_d21\_d32} (Q4) & longitudinal & concatenação de \emph{níveis} T1, D21 e D32 (os mesmos descritores em $T_1,T_2,T_3$) \\
			\bottomrule
		\end{tabular}
	\end{table}

	\begin{table}[htbp]
		\centering
		\caption{Atributos por família \emph{antes} da selecção (L$+$R; denylist de definição, não lista permitida de nomes).
			Q4 $=3\times$ o baseline: cada sufixo em T1, D21 e D32.}
		\label{tab:dims}
		\begin{tabular}{llrr}
			\toprule
			Família & Entra (denylist) & \code{t1\_only} & Q4 \\
			\midrule
			vol & \code{gm/wm/csf\_norm}, \code{MeshVolume} (sai mm$^3$ cru) & 8 & 24 \\
			shape & 12 \code{original\_shape\_*} (sai Mesh/VoxelVolume) & 24 & 72 \\
			texture & 24 \code{original\_glcm\_*} (sai GLRLM/GLSZM/GLDM/NGTDM) & 48 & 144 \\
			firstorder & 16 \code{original\_firstorder\_*} (sai Energy, TotalEnergy) & 32 & 96 \\
			disp & 15 momentos \code{mag\_}/\code{logjac\_}/\code{strain\_fro\_} (sai $u_{x,y,z}$, \code{\_n}, percentis) & 30 & 90 \\
			\bottomrule
		\end{tabular}
	\end{table}

	Sem denylist, só o bloco IBSI dá $107\times 2=214$ colunas em \code{t1\_only} e $642$ em Q4; $n\in[120,231]$.
	A Tabela~\ref{tab:dims} é o espaço que \code{l1\_stable} vê.
	Ainda $p\sim n$ (ou $p>n$) em várias células; o L1 continua necessário.
	Não há modalidade \code{all} (concatenação precoce das cinco famílias) no contraste oficial: esse desenho reproduz $p\gg n$ e mistura papéis no mesmo SVM.

	\subsection{Desenho experimental}
	\label{sec:experiments}
	
	O desfecho é sMCI versus pMCI nas quatro células do fatorial janela clínica $\times$ intervalo (Tabela~\ref{tab:coortes}): \code{36m\_6m} $125/106$, \code{36m\_12m} $121/33$, \code{48m\_6m} $73/120$, \code{48m\_12m} $72/48$.
	Cada conjunto exige três T1 com hipocampo bilateral em $T_1,T_2,T_3$, intervalo na faixa $\pm 2$ meses e desfecho clínico na janela correspondente (Sec.~\ref{subsec:longitudinal-window}).
	Inclusão ampliada de pMCI ($i_3=$ primeiro AD): 46 nos intervalos de 6 meses e 33 nos de 12 meses.
	Não há hold-out fixo nem \emph{downsampling} de classe.
	A generalização avalia-se só no paciente: as três visitas do mesmo sujeito ficam juntas no treino ou no teste.
	O $n$ é o preço de três preditoras elegíveis, intervalo clínico observado e QC; reduz a amostra e impede follow-up incompleto ou informação pós-desfecho de facilitar a tarefa.
	
	Não existe uma coorte única de 243 participantes (153/90) com gap $\ge 6$ --- esse era o protocolo legado.
	Demografia sobretudo em \code{36m\_6m} ($n=231$).
	Nomes de atributos, clínico, leaky e CN$\times$AD em \code{48m\_12m} ($n=120$).
	A fusão tardia (três especificações) corre nas quatro células.
	
	Distinguem-se operações \emph{fixas por imagem} de operações \emph{dependentes da amostra}.
	Pré-processamento, parcelamento, registro e extração correm uma vez por volume, com algoritmos, referências e software fixos; não usam o rótulo sMCI/pMCI nem parâmetros estimados ao agrupar as duas classes.
	Filtros de correlação/variância, $z$-score, selecção, hiperparâmetros e limiar de Youden ajustam-se só no treino externo e aplicam-se ao teste da mesma dobra.
	Limites de QC e templates CN calculam-se uma vez antes da CV.
	
	\subsubsection{Métrica primária e estimadores de AUC}
	\label{subsubsec:evaluation-metrics}
	
	Métrica primária: AUC no paciente sobre probabilidades OOF médias por sujeito.
	Em cada repetição $r\in\{1,\ldots,R\}$ ($R=10$), os cinco testes externos disjuntos cobrem cada participante exatamente uma vez e devolvem $s_{ir}=\widehat{\Pr}(y=1\mid x_i)$ via \code{predict\_proba} (SVM: Platt).
	Agrega-se a média de todos os scores OOF do sujeito (fold$\times$repeat); com folds disjuntos isso coincide com a média das $R=10$ probabilidades.
	Seja $\mathrm{AUC}(\{(y_k,t_k)\})$ a AUC de Mann--Whitney (\code{roc\_auc\_score}; empates a $1/2$).
	\begin{equation}
		\widehat{\mathrm{AUC}}^{(r)}
		=
		\mathrm{AUC}\!\left(\{(y_i,s_{ir})\}_{i=1}^{n}\right),
		\qquad
		\overline{\mathrm{AUC}}
		=
		\frac{1}{R}\sum_{r=1}^{R}\widehat{\mathrm{AUC}}^{(r)}.
		\label{eq:auc-mean-rep}
	\end{equation}
	O estimador primário média primeiro as probabilidades calibradas,
	\begin{equation}
		\bar{s}_i
		=
		\frac{1}{R}\sum_{r=1}^{R}s_{ir},
		\qquad
		\widehat{\mathrm{AUC}}_{\mathrm{agg}}
		=
		\mathrm{AUC}\!\left(\{(y_i,\bar{s}_i)\}_{i=1}^{n}\right),
		\label{eq:auc-agg}
	\end{equation}
	ou, com $s_i\equiv\bar{s}_i$, $n_+=\sum_i y_i$ e $n_-=n-n_+$,
	\begin{equation}
		\mathrm{AUC}_{\mathrm{patient}}
		\equiv
		\widehat{\mathrm{AUC}}_{\mathrm{agg}}
		= \frac{1}{n_+ n_-}
		\sum_{i:\,y_i=1} \sum_{j:\,y_j=0}
		\left(
		\mathbb{I}(s_i > s_j)
		+ \tfrac{1}{2}\,\mathbb{I}(s_i = s_j)
		\right).
		\label{eq:auc-pooled}
	\end{equation}
	Quatro resumos distintos: (i)~AUC de \emph{uma} dobra externa (média$\pm$SD das 50 descreve variação entre modelos, \textbf{não} é incerteza de $\widehat{\mathrm{AUC}}_{\mathrm{agg}}$);
	(ii)~$\widehat{\mathrm{AUC}}^{(r)}$ da repetição;
	(iii)~$\widehat{\mathrm{AUC}}_{\mathrm{agg}}$, o primário;
	(iv)~AUC \emph{pooled}: concatena as $R$ previsões por sujeito sem as mediar (sensibilidade).
	O SD nas tabelas de ponto é o desvio bootstrap de $\widehat{\mathrm{AUC}}_{\mathrm{agg}}$ ($2000$ reamostragens na compilação das planilhas), não o SD das 50 dobras.
	IC95\% e testes de $\Delta$AUC usam bootstrap pareado nos $\{\bar{s}_i\}$ ($B_{\mathrm{boot}}=5000$).
	
	Métricas com limiar usam Youden no treino externo,
	\begin{equation}
		J(\tau) = \mathrm{TPR}(\tau) + \mathrm{TNR}(\tau) - 1
		= \mathrm{TPR}(\tau)-\mathrm{FPR}(\tau),
		\qquad
		\tau^{*} = \arg\max_{\tau} J(\tau).
		\label{eq:youden}
	\end{equation}
	Os classificadores devolvem \code{predict\_proba} (SVM: Platt com \code{probability=True}); não se limiariza \code{decision\_function}.
	$\tau^{*}$ é o limiar da curva ROC (\code{roc\_curve}) que maximiza $J$ nas OOF \code{predict\_proba} do treino externo (\code{cross\_val\_predict} no CV interno) e aplica-se só ao teste correspondente.
	Acurácia, sensibilidade, especificidade e F1: média$\pm$SD nas 50 avaliações.
	A AUC primária não usa limiar.
	
	Constantes: semente $42$ (repetição $r$: $42+1000r$; inner-fold: $42+1000r+\mathrm{fold}$); $k_{\mathrm{out}}=k_{\mathrm{in}}=5$; $|\rho|>0.85$; variância $<0.01$; $\tau=70\%$ em $\pi_g$ (grupo, não coluna); $B=50$ bootstraps; $C=0.1$ na L1 de rastreio; \code{stable-pool-min-timepoints}$=0$.
	(\textbf{Não} $\theta=0.50$.)
	
	\subsubsection{Classificadores e análises comparativas}
	
	Todos os contrastes reportados usam o \textbf{mesmo} SVM~\cite{cortes-1995}.
	O objecto do estudo é o biomarcador (família $\times$ representação temporal $\times$ união tardia), não uma arquitectura nova de aprendizagem.
	Fixar o classificador isola esses contrastes de vieses algorítmicos (capacidade, regularização implícita, importância de atributos).
	O SVM é o padrão em discriminação ADNI com $p$ da ordem de $n$ e $n$ moderado~\cite{cuingnet-2011}: margem máxima, kernels linear ou RBF, resistente a $p\approx n$.
	Na célula mais apertada (\code{36m\_12m}, 33 pMCI) um ensemble não linear decoraria ruído de treino; não se varre RF, elastic-net nem redes.
	O modo \code{l1\_stable} corre \emph{a montante} do SVM e não é o stability selection clássico (sem PFER; Sec.~\ref{sec:feature-selection}).
	
	Kernel linear ou RBF escolhido no CV interno.
	Optuna TPE, 10 ensaios, semente da dobra, maximiza a AUC média interna --- não grid search.
	$C\in[10^{-2},10]$ (log), $\gamma=\mathrm{scale}$, \code{class\_weight}$\in\{\mathrm{None},\mathrm{balanced}\}$, \code{max\_iter}$=2000$.
	Calibração Platt reajustada em cada fit de treino (OOF do treino externo e fit final no treino completo); o teste externo não entra na calibração.
	
	Comparações oficiais: (i)~Q4 versus \code{t1\_only} nas cinco famílias, nas quatro coortes;
	(ii)~fusão tardia em três especificações pré-especificadas (tudo $t_0$; tudo Q4; âncora forma@$t_0$ $\cup$ restante Q4).
	Extra em \code{48m\_12m}, mesmo SVM: clínico de $t_0$; clínica$+$forma \code{t1\_only}; CN$\times$AD (sanidade); leaky (apêndice).
	A fusão precoce entre famílias de imagem (\code{--modality all}) \textbf{não} entra: concatena classes completas e reproduz $p\gg n$.
	
	Inferência sobre $\{\bar{s}_i\}$.
	Teste contra o acaso: permutação de rótulos $B_{\mathrm{perm}}=5000$ com scores congelados --- a pipeline nested \textbf{não} se reexecuta.
	$p=(r+1)/(B_{\mathrm{perm}}+1)$.
	Contrastes pareados (Q4 vs \code{t1\_only}; late vs teto unimodal): bootstrap de participantes $B_{\mathrm{boot}}=5000$, IC percentílico 95\%, mesma correcção $(r+1)/(B+1)$.
	As duas representações partilham o protocolo $10\times 5$; o emparelhamento é nos scores agregados por paciente, não nas AUCs de dobra.
	No contraste temporal unimodal, o $p$ (quando reportado) é \textbf{unilateral} (H$_1$: $\mathrm{AUC}(\mathrm{Q4})>\mathrm{AUC}(\code{t1\_only})$); o $q$ é FDR de Benjamini--Hochberg nas \textbf{cinco} famílias da mesma coorte.
	$p$ bilateral reserva-se a contrastes sem direcção pré-especificada (p.ex.\ imagem vs clínico).
	Os pontos das Tabelas~\ref{tab:uni_4}--\ref{tab:late_4} são $\widehat{\mathrm{AUC}}_{\mathrm{agg}}$; testes pareados do $\Delta$ ficam para a compilação estatística sobre estes CSVs e não se herdaram do protocolo \code{t1\_deltas}.
	
	\subsection{Validação aninhada no paciente}
	\label{sec:modelagem}
	
	CV aninhada no identificador do paciente (\code{StratifiedKFold}).
	A matriz de classificação é uma linha por paciente (colunas T1/D21/D32 já no wide da representação); as dobras partem-se por paciente, não por imagem: as três visitas do mesmo $i$ ficam todas no treino ou todas no teste.
	Loop externo: $k_{\mathrm{out}}=5$ folds $\times$ $R=10$ repetições $=50$ avaliações por configuração.
	Semente da repetição $r$: $42+1000r$; o \code{StratifiedKFold} externo usa essa semente.
	Loop interno, em cada dobra externa: $k_{\mathrm{in}}=5$, semente $42+1000r+\mathrm{fold}$.
	
	Por dobra externa, nesta ordem, só no treino:
	\begin{enumerate}
		\item Pool \code{l1\_stable}: constrói-se \textbf{uma vez} neste treino externo ($B=50$; Sec.~\ref{sec:feature-selection}); \textbf{não} se reconstrói em cada split interno.
		\item Pipeline sklearn: \code{CorrVarSelector} (correlação $|\rho|>0.85$ depois variância $<0.01$, restrito ao pool) $\to$ \code{StandardScaler} $\to$ SVM (\code{probability=True}).
		\item Optuna TPE, 10 ensaios, maximiza a AUC média das cinco inner-folds (scores \code{predict\_proba}).
		\item $\tau^{*}$ (Eq.~\ref{eq:youden}) estima-se com \code{cross\_val\_predict} no treino externo (mesmo $k_{\mathrm{in}}$) e aplica-se só ao teste desta dobra.
		\item A pipeline vencedora refaz-se no treino externo completo (filtros, $z$-score e Platt reajustados); o teste recebe só \code{transform}/\code{predict\_proba}.
	\end{enumerate}
	No split interno, \code{CorrVarSelector} e \code{StandardScaler} \textbf{reajustam-se} no inner-train, restritos ao pool da dobra; o L1 de rastreio não corre de novo.
	Se os filtros internos zerarem o espaço, o \code{CorrVarSelector} devolve o pool (não a classe inteira).
	A nested cobre o modelamento sample-dependent (pool \code{l1\_stable}, filtros, $z$-score, hiperparâmetros, Youden, Platt).
	Não cobre pré-processamento de imagem, QC nem templates CN --- esses correm uma vez antes da CV.
	
%	\begin{figure}[htbp]
%		\centering
%		\includegraphics[width=1.0\textwidth]{graphics/nestedcv.pdf}
%		\caption{CV aninhada no paciente e fronteira de vazamento.
%			Dobras externas estratificadas ($5\times 10$) por identificador.
%			Abaixo da linha: pool \code{l1\_stable} por grupo ($\pi_g\ge 70\%$), \code{CorrVarSelector}+$z$-score, Optuna TPE (10 ensaios) e Youden --- só treino externo.
%			Ranking primário: AUC paciente nas OOF médias por sujeito; a média das AUCs de dobra é descritiva.}
%		\label{fig:nestedcv}
%	\end{figure}
	
	\subsection{Selecção de atributos (\code{l1\_stable})}
	\label{sec:feature-selection}

	Volume prognóstico: fracções teciduais \emph{dentro da ROI} e \code{original\_shape\_MeshVolume} já /ICV.
	Deslocamento não leva ICV.
	Tabelas long $\to$ uma linha por paciente; só hipocampo L/R.
	Representações oficiais: \code{t1\_only} e Q4 (\code{t1\_d21\_d32}).
	A chave \code{anatomical\_key} colapsa o instante: ROI $+$ lado $+$ basename (T1/D21/D32 retirados). Exemplo: \code{hippocampus\_L\_T1\_gm\_norm} e \code{hippocampus\_L\_D21\_gm\_norm} partilham a mesma chave.

	Todos os números reportados usam o modo \code{l1\_stable}.
	É um filtro descritivo de frequência $\ell_1$ por bootstrap~\cite{bach-2008,meinshausen-2010}: \textbf{não} é o stability selection clássico e \textbf{não} herda PFER.
	Por dobra externa o modo corre \textbf{uma vez} no outer-train e devolve um \emph{pool} de colunas; o SVM nunca vê a classe completa.
	O CV interno \textbf{não} reconstrói este pool --- só reajusta \code{CorrVarSelector} e \code{StandardScaler} restritos a ele.

	Cadeia no outer-train, pela ordem:
	\begin{enumerate}
		\item $z$-score de \textbf{todas} as colunas no outer-train (média e DP só deste treino).
		\item Correlação $|\rho|>0.85$ \emph{uma vez}, greedy: em cada par colinear fica a coluna de maior $|\mathrm{corr}(x,y)|$ com o rótulo $y$.
		Pares com o mesmo \code{anatomical\_key} (T1 vs D21 vs D32 do mesmo descritor) \textbf{nunca} se removem nesta poda --- o filtro não trata tempos do mesmo biomarcador como mecanismos rivais.
		\item Limiar de variância $<0.01$ nas colunas já $z$-score; se o filtro zerar o espaço, retêm-se colunas com $\mathrm{var}>0$ (\textbf{não} se restaura a classe).
		\item $B=50$ amostras com reposição sobre o espaço já podado.
		Réplica com uma só classe: descartada.
		Cada réplica biclasse: $z$-score \emph{da própria amostra} e logística $\ell_1$
	\begin{equation}
		\begin{aligned}
			\hat{\boldsymbol{\beta}}^{(b)},\,\hat{\beta}_0^{(b)}
			&= \arg\min_{\boldsymbol{\beta},\beta_0}
			\left\{
			\sum_{i=1}^{n_{\mathrm{train}}^{(b)}}
			\ell\!\left(y_i,\,\sigma(\mathbf{x}_i^{\top}\boldsymbol{\beta}+\beta_0)\right)
			+ C^{-1}\|\boldsymbol{\beta}\|_1
			\right\}, \\
			\ell(y,p) &= -y\log p-(1-y)\log(1-p),
		\end{aligned}
		\label{eq:l1-logistic}
	\end{equation}
		solver SAGA, $C=0.1$ fixo (sklearn: $C$ inverso da força; intercepto $\beta_0$ sem penalidade $\ell_1$).
		$\mathcal{M}^{(b)}=\{j:\lvert\hat{\beta}^{(b)}_j\rvert>10^{-9}\}$.
		Réplica com $\mathcal{M}^{(b)}=\emptyset$: descartada (não entra em $B_{\mathrm{eff}}$; não é keep-filters).
	\end{enumerate}
	A frequência é \textbf{por grupo}, não por coluna.
	$B_{\mathrm{eff}}\le 50$ conta só réplicas biclasse com $\mathcal{M}^{(b)}\ne\emptyset$.
	O grupo $g$ conta na réplica $b$ se \emph{alguma} coluna de $g$ está em $\mathcal{M}^{(b)}$:
	\begin{equation}
		\pi_g = \frac{1}{B_{\mathrm{eff}}}\sum_{b=1}^{B_{\mathrm{eff}}} \mathbb{I}\!\left(g\cap\mathcal{M}^{(b)}\ne\emptyset\right).
		\label{eq:stability-freq}
	\end{equation}
	Com \code{stable-pool-min-timepoints}$=0$, entra no pool todo $g$ com $\pi_g\ge\tau=0.70$.
	Se $g$ passa, o pool inclui \textbf{todas} as colunas originais desse \code{anatomical\_key} (T1, D21 e D32 quando existem), mesmo as que a L1 não escolheu e as que corr/var tinham cortado antes do bootstrap.
	Por isso Q4 de primeira ordem tende a $n=3$ por descritor aceite: o grupo entra e os três instantes voltam juntos.
	
	Se nenhum grupo passa (inclui $B_{\mathrm{eff}}=0$): \emph{corr\_fallback} --- nomes que sobreviveram a corr+var no outer-train $z$-score, nunca a classe inteira.
	Se esse conjunto também é vazio: \emph{var\_fallback} --- uma coluna de máxima variância no outer-train bruto.
	Entre colineares de chaves distintas \code{l1\_stable} escolhe um representante, não o ``melhor Haralick''.
	
	Depois do pool, a pipeline sklearn \textbf{repete} corr ($|\rho|>0.85$) depois var ($<0.01$), agora nos valores \emph{brutos} do conjunto em que o \code{fit} corre (inner-train no Optuna; outer-train no fit final), restrita ao pool; em seguida o \code{StandardScaler}.
	Se essa segunda poda zerar o espaço, devolve-se o pool, não a classe.
	O CV interno não reconstrói o bootstrap L1.

	\subsection{Protocolo de modelagem}
	\label{sec:protocolo-oficial}

	Protocolo oficial: SVM com \code{probability=True} (Platt), \code{l1\_stable}, nested estratificada $5$ folds $\times$ $10$ repetições (semente $42+1000r$), 50 avaliações externas.
	Dobras por identificador de paciente.

	\paragraph{Hiperparâmetros.}
	Optuna TPE, 10 ensaios, AUC média do CV interno.
	SVM: $C\in[10^{-2},10]$ (log), kernel $\in\{\mathrm{linear},\mathrm{RBF}\}$, $\gamma=\mathrm{scale}$, \code{class\_weight} $\in\{\mathrm{None},\mathrm{balanced}\}$, \code{max\_iter}$=2000$.
	Um único classificador em todos os contrastes; não se varre RF nem elastic-net.

	\paragraph{Youden.}
	$\tau^{*}$ maximiza $J=\mathrm{TPR}+\mathrm{TNR}-1$ nas OOF do treino externo e aplica-se só ao teste dessa dobra.
	A AUC primária não usa limiar.

	\paragraph{AUC no paciente.}
	Média de todos os scores OOF do sujeito (fold$\times$repeat); com $5$ testes disjuntos isso é a média das $R=10$ probabilidades.
	Uma AUC de Mann--Whitney nos $n$ pares $(y_i,\bar{s}_i)$ (Eqs.~\ref{eq:auc-agg}--\ref{eq:auc-pooled}).
	O $\pm$ nas tabelas de ponto é o SD bootstrap dessa AUC paciente ($2000$ reamostragens), \textbf{não} o DP das 50 dobras.

	\subsection{Fusão multimodal e preditores clínicos}
	\label{sec:fusao}

	A fusão tardia \textbf{não} retreina um meta-classificador.
	Reutiliza os scores OOF \code{predict\_proba} dos cinco SVMs unimodais já ajustados no mesmo nested (mesma semente $\Rightarrow$ mesmos splits).
	Por (repetição $r$, dobra $f$), alinhamento inner-join em \code{ID\_PT} e média simples
	\begin{equation}
		s_{i,r,f}^{\mathrm{late}}
		=
		\frac{1}{5}\sum_{k=1}^{5} s_{i,r,f}^{(k)}.
		\label{eq:late-mean}
	\end{equation}
	A AUC paciente (Eqs.~\ref{eq:auc-agg}--\ref{eq:auc-pooled}) calcula-se sobre $\{s_{i,r,f}^{\mathrm{late}}\}$ como nas famílias isoladas.
	Pesos iguais; sem stacking, sem SVM de segundo nível, sem \code{weighted}.
	Cada ramo tem o seu pool \code{l1\_stable}, scaler, Optuna e Platt; forma não vê 144 GLCM e $p\gg n$ da concatenação não se aplica.
	Três especificações \textbf{pré-especificadas} (não escolhidas a posteriori):
	tudo \code{t1\_only}; tudo Q4; âncora \textbf{forma em \code{t1\_only} $\cup$ volume, textura, deslocamento e primeira ordem em Q4} (cinco ramos, não seis: shape não entra também em Q4).
	A AUC primária da fusão não usa limiar.
	Sens/Esp da linha late, se reportadas, herdam o $\tau^{*}$ do primeiro ramo da spec (não se reestima Youden no score médio).
	A fusão precoce entre famílias de imagem está fora do paper.

	Preditores clínicos de $t_0$ (extra, mesmo SVM): sexo, idade, MMSE, ADAS e FAQ (sem CDR global).
	A fusão clínica$+$imagem concatena, no treino externo, esses cinco escores $z$-score à forma em \code{t1\_only}; um único SVM.
	O protocolo \emph{leaky} (apêndice) estima $z$-score, filtros e pool \code{l1\_stable} no conjunto completo antes da CV.

	\subsection{Métrica primária e papéis do bootstrap}
\label{sec:stats}


A métrica oficial é a AUC no paciente $\widehat{\mathrm{AUC}}_{\mathrm{agg}}$ (Eqs.~\ref{eq:auc-agg}--\ref{eq:auc-pooled}), não a média das 50 AUCs de fold nem a AUC \emph{pooled}.
CV aninhada $5\times 10$ no identificador do paciente: em cada repetição $r$ o sujeito $i$ cai uma vez no teste e recebe a probabilidade OOF $s_{ir}$ (\code{predict\_proba}; SVM com Platt).
Primeiro tira-se a média,
\[
	\bar{s}_i = \frac{1}{10}\sum_{r=1}^{10} s_{ir},
\]
e calcula-se uma única vez a AUC de Mann--Whitney nos $n$ pares $(y_i,\bar{s}_i)$.
A AUC primária não usa limiar de Youden.

Há \textbf{três} bootstraps, com papéis separados.
Nenhum reexecuta a pipeline nested.
A permutação de rótulos \textbf{não} é bootstrap.

\paragraph{A.~\code{l1\_stable} (dentro da CV).}
Modo \code{l1\_stable}, uma vez por treino externo: $z$-score $\to$ corr $0.85$ $\to$ var $0.01$ $\to$ $B=50$ L1 (Eqs.~\ref{eq:l1-logistic}--\ref{eq:stability-freq}).
Frequência $\pi_g$ por \code{anatomical\_key}; entra o grupo com $\pi_g\ge 70\%$ e \textbf{todos} os seus instantes.
Pool vazio $\to$ \code{corr\_fallback}, nunca a classe inteira.
Esse bootstrap \textbf{não} estima AUC: só escolhe colunas.
Não é o stability selection clássico e não herda PFER~\cite{bach-2008,meinshausen-2010}.
Detalhe: Sec.~\ref{sec:feature-selection}.

\paragraph{B.~Desvio nas tabelas e barras de erro (depois da CV).}
Scores congelados $\{\bar{s}_i\}$.
$B_{\mathrm{sd}}=2000$ reamostragens de pacientes com reposição; em cada uma recalcula-se $\widehat{\mathrm{AUC}}_{\mathrm{agg}}$.
O $\pm$ reportado (Tabelas~\ref{tab:uni_4}--\ref{tab:late_4}) é o desvio-padrão dessas $2000$ AUCs,
\[
	\widehat{\sigma}_{\mathrm{boot}}
	= \mathrm{sd}\bigl\{\widehat{\mathrm{AUC}}_{\mathrm{agg}}^{(b)}\bigr\}_{b=1}^{B_{\mathrm{sd}}},
\]
\textbf{não} o DP das 50 folds e \textbf{não} um intervalo de confiança a 95\%.

\paragraph{C.~Inferência do $\Delta$AUC (depois da CV).}
Contrastes pareados no mesmo sujeito (p.ex.\ Q4 vs \code{t1\_only}): $B_{\mathrm{boot}}=5000$ reamostragens,
\[
	\Delta^{(b)}
	= \widehat{\mathrm{AUC}}_{\mathrm{agg}}(A^{(b)})
	- \widehat{\mathrm{AUC}}_{\mathrm{agg}}(B^{(b)}).
\]
IC95\% percentílico $[2{,}5,\,97{,}5]$ da distribuição empírica de $\Delta^{(b)}$.
\[
	p_{\mathrm{unilateral}}
	= \frac{\#\{\Delta^{(b)}\le 0\}+1}{B_{\mathrm{eff}}+1}
	\qquad
	(\mathrm{H}_1{:}\ \mathrm{AUC}_A>\mathrm{AUC}_B).
\]
$p$ bilateral $=2\min(p_{\mathrm{uni}},1-p_{\mathrm{uni}})$, reservado a contrastes sem direção pré-especificada (imagem vs clínico).
No contraste temporal unimodal, o $q$ é FDR de Benjamini--Hochberg nos cinco $p$ unilaterais da mesma coorte.

\paragraph{Permutação contra o acaso.}
Embaralha-se $y$ com $\{\bar{s}_i\}$ fixos, $B_{\mathrm{perm}}=5000$;
$p=(r+1)/(B_{\mathrm{perm}}+1)$, em que $r$ é o número de permutações com AUC $\ge$ à observada.

$\alpha=0.05$.
O contraste temporal principal é Q4 versus \code{t1\_only} nas cinco famílias, nas quatro coortes.
Pontos das tabelas de resultados são $\widehat{\mathrm{AUC}}_{\mathrm{agg}}$; $p$/$q$ pareados destes CSVs não se herdaram do protocolo \code{t1\_deltas}.

	\subsection{Testes estatísticos}
	\label{sec:testes}

	Todos os contrastes de AUC usam os scores agregados por paciente, não as AUCs de dobra.
	\begin{itemize}
		\item Permutação de rótulos com scores congelados ($B_{\mathrm{perm}}=5000$): $p=(r+1)/(B+1)$; o pipeline nested \textbf{não} se reexecuta.
		\item Bootstrap pareado do $\Delta$AUC ($B_{\mathrm{boot}}=5000$), IC percentílico a 95\%, $p$ unilateral (H$_1$: AUC$_A>$AUC$_B$) e bilateral.
		\item FDR de Benjamini--Hochberg sobre os $p$ unilaterais das cinco famílias, dentro de cada coorte.
		\item Demografia em \code{36m\_6m}: Mann--Whitney na idade, $\chi^2$ no sexo; classificadores univariados de idade e de sexo com permutação.
	\end{itemize}
	$\alpha=0.05$.
	
	% =============================================================================
	\section{Resultados}
	\label{sec:resultados}

	Apresentação: \textbf{uma matriz para as quatro células}, não quatro relatos paralelos.
	Protocolo fixo SVM $+$ \code{l1\_stable}.
	Métrica: AUC paciente $\widehat{\mathrm{AUC}}_{\mathrm{agg}}$; $\pm$ é SD bootstrap ($B=2000$).
	$\Delta=\mathrm{AUC}(\mathrm{Q4})-\mathrm{AUC}(\code{t1\_only})$.
	Claim multimodal e nomes de atributos: \code{48m\_12m} ($n=120$); o fatorial testa se o padrão se replica.

	\subsection{Demografia}
	\label{subsec:demo}

	Em \code{36m\_6m} ($n=231$), idade e sexo não diferem entre sMCI e pMCI (Mann--Whitney $p=0.071$; $\chi^2$ $p=0.113$).
	Classificador só com idade: AUC $0.560$; só com sexo: $0.452$.
	A classificação hipocampal não se reduz a demografia (Tabela~\ref{tab:clin_cohorts}).

	\subsection{Uniclasse: Q4 versus \texorpdfstring{\code{t1\_only}}{t1\_only} nas quatro coortes}
	\label{subsec:unimodal}

	A Tabela~\ref{tab:uni_4} é o resultado principal unimodal.

	\begin{sidewaystable}[htbp]
		\centering
		\footnotesize
		\setlength{\tabcolsep}{3.5pt}
		\caption{AUC no paciente, SVM.
			Em cada célula: \code{t1\_only} / Q4 / $\Delta$.
			Negrito: maior das duas representações naquela família$\times$coorte.
			$n$: 231, 154, 193, 120.}
		\label{tab:uni_4}
		\begin{tabular}{@{}l ccc ccc ccc ccc@{}}
			\toprule
			& \multicolumn{3}{c}{\code{36m\_6m}} & \multicolumn{3}{c}{\code{36m\_12m}} & \multicolumn{3}{c}{\code{48m\_6m}} & \multicolumn{3}{c}{\code{48m\_12m}} \\
			\cmidrule(lr){2-4}\cmidrule(lr){5-7}\cmidrule(lr){8-10}\cmidrule(lr){11-13}
			Família & $t_0$ & Q4 & $\Delta$ & $t_0$ & Q4 & $\Delta$ & $t_0$ & Q4 & $\Delta$ & $t_0$ & Q4 & $\Delta$ \\
			\midrule
			vol & 0.738 & \textbf{0.752} & $+0.014$ & \textbf{0.707} & 0.687 & $-0.019$ & 0.714 & \textbf{0.759} & $+0.045$ & 0.632 & \textbf{0.662} & $+0.030$ \\
			 & $\pm$0.032 & $\pm$0.032 & & $\pm$0.047 & $\pm$0.052 & & $\pm$0.038 & $\pm$0.036 & & $\pm$0.054 & $\pm$0.051 & \\
			shape & 0.789 & \textbf{0.795} & $+0.006$ & \textbf{0.739} & 0.732 & $-0.006$ & \textbf{0.801} & 0.797 & $-0.004$ & \textbf{0.761} & 0.755 & $-0.006$ \\
			 & $\pm$0.030 & $\pm$0.030 & & $\pm$0.046 & $\pm$0.047 & & $\pm$0.035 & $\pm$0.035 & & $\pm$0.043 & $\pm$0.043 & \\
			texture & 0.727 & \textbf{0.739} & $+0.012$ & 0.724 & \textbf{0.745} & $+0.021$ & 0.740 & \textbf{0.759} & $+0.019$ & 0.711 & \textbf{0.762} & $+0.052$ \\
			 & $\pm$0.034 & $\pm$0.033 & & $\pm$0.051 & $\pm$0.051 & & $\pm$0.037 & $\pm$0.036 & & $\pm$0.050 & $\pm$0.044 & \\
			firstorder & 0.752 & \textbf{0.768} & $+0.016$ & 0.741 & \textbf{0.760} & $+0.020$ & 0.773 & \textbf{0.797} & $+0.024$ & 0.719 & \textbf{0.737} & $+0.019$ \\
			 & $\pm$0.032 & $\pm$0.031 & & $\pm$0.051 & $\pm$0.050 & & $\pm$0.035 & $\pm$0.033 & & $\pm$0.048 & $\pm$0.048 & \\
			disp & \textbf{0.628} & 0.501 & $-0.127$ & 0.476 & \textbf{0.492} & $+0.016$ & \textbf{0.609} & 0.571 & $-0.038$ & 0.524 & \textbf{0.621} & $+0.097$ \\
			 & $\pm$0.036 & $\pm$0.039 & & $\pm$0.057 & $\pm$0.056 & & $\pm$0.043 & $\pm$0.043 & & $\pm$0.055 & $\pm$0.052 & \\
			\bottomrule
		\end{tabular}
	\end{sidewaystable}

	\paragraph{Leitura da matriz.}
	\begin{itemize}
		\item \textbf{Textura e primeira ordem:} $\Delta>0$ nas quatro células. Máximo de textura $+0.052$ em \code{48m\_12m} ($0.711\to 0.762$).
		\item \textbf{Volume:} $\Delta>0$ em três células; quebra em \code{36m\_12m} ($-0.019$).
		\item \textbf{Forma:} teto estático. Q4 nunca ganha de forma material ($|\Delta|\le 0.006$). Melhor unimodal em \code{36m\_6m}/\code{48m\_6m}/\code{48m\_12m} é shape-$t_0$ (ou empate com textura Q4 no claim: $0.762$ vs $0.761$).
		\item \textbf{Deslocamento:} fraco. Em intervalos de 12 meses o pool \code{l1\_stable} esvazia-se na maioria das dobras (\code{corr\_fallback}); o $\Delta$ positivo em \code{48m\_12m} \textbf{não} se interpreta como selecção estável.
	\end{itemize}
	Q4 agrega nas famílias de intensidade/volume e \textbf{não} passa o teto shape-$t_0$ de forma consistente.

	\subsection{Quantos e quais atributos}
	\label{subsec:atributos-sel}

	A Tabela~\ref{tab:nfeat} é a média de colunas que o SVM vê após \code{l1\_stable} ($50$ avaliações).
	A Tabela~\ref{tab:feat_claim} lista os grupos com cobertura $\ge 20\%$ em \code{48m\_12m}.

	\begin{table}[htbp]
		\centering
		\footnotesize
		\caption{Média de atributos seleccionados (min--máx nas 50 dobras).
			Entra: Tabela~\ref{tab:dims}.
			Disp em 12\,meses: maioria \code{corr\_fallback}.}
		\label{tab:nfeat}
		\begin{tabular}{ll rr rr rr rr}
			\toprule
			& & \multicolumn{2}{c}{\code{36m\_6m}} & \multicolumn{2}{c}{\code{36m\_12m}} & \multicolumn{2}{c}{\code{48m\_6m}} & \multicolumn{2}{c}{\code{48m\_12m}} \\
			\cmidrule(lr){3-4}\cmidrule(lr){5-6}\cmidrule(lr){7-8}\cmidrule(lr){9-10}
			Fam. & enc & média & faixa & média & faixa & média & faixa & média & faixa \\
			\midrule
			vol & $t_0$ & 1.00 & 1--1 & 1.18 & 1--4 & 1.08 & 1--2 & 2.14 & 1--5 \\
			vol & Q4 & 6.28 & 5--11 & 4.98 & 3--12 & 6.14 & 3--9 & 5.16 & 3--13 \\
			shape & $t_0$ & 1.64 & 1--2 & 1.76 & 1--12 & 2.00 & 1--3 & 1.52 & 1--3 \\
			shape & Q4 & 7.80 & 3--12 & 4.08 & 3--27 & 7.14 & 6--12 & 5.12 & 3--9 \\
			texture & $t_0$ & 1.66 & 1--4 & 2.22 & 1--8 & 1.52 & 1--3 & 1.62 & 1--10 \\
			texture & Q4 & 5.38 & 3--15 & 3.74 & 3--25 & 5.32 & 3--12 & 3.60 & 3--6 \\
			firstorder & $t_0$ & 3.00 & 1--4 & 1.62 & 1--7 & 2.84 & 2--4 & 1.00 & 1--1 \\
			firstorder & Q4 & 10.86 & 4--15 & 3.82 & 3--29 & 9.92 & 6--15 & 3.00 & 3--3 \\
			disp & $t_0$ & 1.50 & 1--2 & 13.90 & 1--18 & 1.42 & 1--2 & 15.74 & 1--18 \\
			disp & Q4 & 4.92 & 1--10 & 33.08 & 1--50 & 3.72 & 1--9 & 41.62 & 1--50 \\
			\bottomrule
		\end{tabular}
	\end{table}

	\begin{table}[htbp]
		\centering
		\footnotesize
		\caption{Grupos retidos em $\ge 20\%$ das 50 avaliações, \code{48m\_12m}.
			Cobertura = fracção de folds com o biomarcador (qualquer visita).
			Em Q4, T1/D21/D32 do mesmo \code{anatomical\_key} tendem a entrar juntos.}
		\label{tab:feat_claim}
		\begin{tabular}{lllrr}
			\toprule
			Família & enc & Grupo & cob.\ (\%) & notas temporais \\
			\midrule
			vol & $t_0$ & R \code{csf\_norm} & 84 & T1 \\
			vol & $t_0$ & L \code{csf\_norm} & 50 & T1 \\
			vol & Q4 & L \code{csf\_norm} & 84 & T1$=$D21$=$D32 \\
			vol & Q4 & R \code{csf\_norm} & 68 & T1$=$D21$=$D32 \\
			shape & $t_0$ & L SurfaceVolumeRatio & 68 & T1 \\
			shape & $t_0$ & R SurfaceVolumeRatio & 64 & T1 \\
			shape & Q4 & R SurfaceVolumeRatio & 82 & $\approx$ três visitas \\
			shape & Q4 & L SurfaceVolumeRatio & 78 & $\approx$ três visitas \\
			texture & $t_0$ & L InverseVariance & 62 & T1 \\
			texture & Q4 & L InverseVariance & 52 & T1$=$D21$=$D32 \\
			firstorder & $t_0$/Q4 & L Entropy / Uniformity & 42 / 38 & Q4: três visitas \\
			disp & $t_0$/Q4 & momentos \code{mag}/\code{logjac}/\code{strain\_fro} & $\sim$84--92 & \code{corr\_fallback} \\
			\bottomrule
		\end{tabular}
	\end{table}

	L1 corta a classe a poucas colunas (tipicamente 1--6), excepto deslocamento em 12\,meses.
	Nomes = frequência, não assinatura única: SurfaceVolumeRatio e \code{csf\_norm} são estáveis; InverseVariance e Entropy/Uniformity são os representantes de textura e primeira ordem.

	\subsection{Multimodal: fusão tardia nas quatro coortes}
	\label{subsec:multimodal}

	A Tabela~\ref{tab:late_4} compara as três especificações ao teto unimodal da mesma célula (máximo da Tabela~\ref{tab:uni_4}).

	\begin{table}[htbp]
		\centering
		\caption{AUC paciente da fusão tardia (média de cinco SVMs) versus teto unimodal.
			Âncora: shape@$t_0$ $\cup$ vol/texture/disp/firstorder em Q4.}
		\label{tab:late_4}
		\begin{tabular}{lrrrr}
			\toprule
			Coorte & Teto uni. & Late $t_0$ & Late Q4 & Âncora \\
			\midrule
			\code{36m\_6m} & 0.795 & $0.808\pm 0.029$ & $0.814\pm 0.028$ & $0.814\pm 0.029$ \\
			\code{36m\_12m} & 0.760 & $0.769\pm 0.045$ & $0.793\pm 0.039$ & $0.792\pm 0.041$ \\
			\code{48m\_6m} & 0.801 & $0.816\pm 0.035$ & $0.828\pm 0.032$ & $\mathbf{0.833\pm 0.032}$ \\
			\code{48m\_12m} & 0.762 & $0.761\pm 0.044$ & $\mathbf{0.802\pm 0.040}$ & $\mathbf{0.802\pm 0.040}$ \\
			\bottomrule
		\end{tabular}
	\end{table}

	Late Q4 e âncora superam o teto unimodal nas quatro células.
	No claim, late $t_0$ \textbf{não} passa o teto ($0.761$ vs shape-$t_0$ $0.761$); quem passa é a união \emph{com} Q4 ($0.802$).
	Âncora $\approx$ late Q4: a forma já saturava no baseline; o ganho é empilhar os papéis que se movem no tempo.

	\subsection{Clínico, sanidade e literatura}
	\label{subsec:controlos}

	Em \code{48m\_12m}, SVM só com sexo, idade, MMSE, ADAS e FAQ em $t_0$: $0.812\pm 0.040$.
	Fusão precoce clínica$+$forma \code{t1\_only}: $0.843\pm 0.036$.
	A RM hipocampal isolada não supera o clínico; acrescentá-la sobe o ponto sem ser o claim das perguntas 1--3.

	CN versus AD no mesmo stack (volume Q4, \code{48m\_12m}): $0.833\pm 0.030$ --- a tarefa fácil funciona; sMCI/pMCI permanece a difícil.
	Leaky (volume Q4, selecção global): $0.660$ versus nested $0.662$.

	A Tabela~\ref{tab:literature_comparison} situa o teto unimodal e a fusão tardia face a estudos sMCI/pMCI representativos.
	Os números \textbf{não} são comparáveis: diferem horizonte, ROI, preditores clínicos e validação.

\begin{sidewaystable}[htbp]
	\centering
	\footnotesize
	\setlength{\tabcolsep}{4pt}
	\renewcommand{\arraystretch}{1.15}
	\caption{Comparação indicativa de desempenho prognóstico publicado em sMCI versus pMCI e dos resultados ADNI deste estudo sob validação aninhada no paciente.}
	\label{tab:literature_comparison}
	\begin{tabular}{@{}lllccc@{}}
		\toprule
		\textbf{Estudo} & \textbf{Base} & \textbf{Atributos} &
		\textbf{Validação} & \textbf{AUC} & \textbf{Sens / Esp (\%)} \\
		\midrule
		\citet{bucholc-2023} & NACC &
		Cognitivo + funcional &
		3 visitas; CV retido &
		$0.601 \pm 0.054$ & $95.4 \pm 1.8$ / $24.7 \pm 12.0$ \\
		\citet{bucholc-2023} & NACC &
		Cognitivo + funcional &
		4 visitas; CV retido &
		$0.756 \pm 0.103$ & $92.9 \pm 4.6$ / $58.3 \pm 19.1$ \\
		\citet{grassi-2019} & ADNI &
		Sociodemográfico + clínico &
		Baseline; CV 10 folds &
		$0.88$ & $77.7$ / $79.9$ \\
		\citet{moradi-2015} & ADNI &
		Densidade de GM (cérebro inteiro) &
		Baseline; CV 10 folds &
		$0.766$ & $88.9$ / $51.6$ \\
		\citet{moradi-2015} & ADNI &
		GM (cérebro inteiro) + cognitivo &
		Baseline; CV 10 folds &
		$0.902$ & $86.7$ / $73.6$ \\
		\citet{ye-2012}$^{\S}$ & ADNI &
		Cérebro inteiro + cognitivo &
		Baseline; stability selection &
		$0.859$ & \textemdash\ / \textemdash \\
		\citet{casanova-2013}$^{\P}$ & ADNI &
		Tecidos (cérebro inteiro) + cognitivo &
		Baseline; CV aninhada &
		\textemdash & $45.8$ / $75.5$ \\
		\citet{guo-2017a}$^{\|}$ & ADNI &
		Córtex &
		Baseline; LOOCV &
		\textemdash & $85.0$ / $84.8$ \\
		\citet{cuingnet-2011} & ADNI &
		Volume hipocampal &
		Baseline; LOOCV &
		\textemdash & $62.0$ / $69.0$ \\
		\midrule
		\textbf{Este trabalho (teto $t_0$)}$^{\dagger}$ & ADNI &
		Forma hipocampal &
		$t_0$; nested CV (SVM) &
		$\mathbf{0.761 \pm 0.043}$ & $\mathbf{70.5 \pm 15.4}$ / $\mathbf{63.7 \pm 12.7}$ \\
		\textbf{Este trabalho (Q4 unimodal)} & ADNI &
		Textura GLCM (níveis T1+D21+D32) &
		3 visitas; nested CV &
		$\mathbf{0.762 \pm 0.044}$ & $\mathbf{68.3 \pm 17.8}$ / $\mathbf{67.3 \pm 15.1}$ \\
		\textbf{Este trabalho (fusão tardia)}$^{\ddagger}$ & ADNI &
		Âncora / late Q4 (5 SVMs) &
		3 visitas; nested CV &
		$\mathbf{0.802 \pm 0.040}$ & $\mathbf{74.2 \pm 18.3}$ / $\mathbf{66.8 \pm 14.8}$ \\
		\textbf{Este trabalho} & ADNI &
		Clínico (sexo, idade, MMSE, ADAS, FAQ) &
		$t_0$; nested CV (SVM) &
		$\mathbf{0.812 \pm 0.040}$ & $\mathbf{70.0 \pm 14.2}$ / $\mathbf{70.9 \pm 14.1}$ \\
		\textbf{Este trabalho} & ADNI &
		Clínico $+$ forma $t_0$ &
		nested CV (SVM) &
		$\mathbf{0.843 \pm 0.036}$ & $\mathbf{77.1 \pm 13.0}$ / $\mathbf{73.2 \pm 13.6}$ \\
		\bottomrule
	\end{tabular}
	\vspace{0.5em}
	\begin{flushleft}
		\scriptsize
		\textit{Nota:}
		Linhas próprias: coorte \code{48m\_12m} ($n=120$).
		$^{\dagger}$Teto unimodal de baseline (forma, \code{t1\_only}).
		$^{\ddagger}$Fusão tardia; Sens/Esp da linha late Q4 (Youden por dobra).
		Antes de $\pm$ na AUC: estimando paciente; depois: SD bootstrap ($2000$).
		Sens/Esp: média $\pm$ SD nas 50 avaliações externas.
	\end{flushleft}
\end{sidewaystable}

	\section{Discussão}
	\label{sec:discussao}

	A pergunta era se níveis longitudinais hipocampais agregam discriminação sMCI/pMCI além do $t_0$.
	A resposta depende do \emph{papel} da família e de uniclasse versus união.

	\paragraph{Uniclasse agrega, não passa o teto.}
	Q4 sobe volume, textura e primeira ordem; a forma já discrimina no $t_0$ e Q4 é redundante (o filtro de correlação nem deixa T1 competir com D21/D32 do mesmo descritor como se fossem mecanismos distintos).
	Nenhuma família Q4 substitui de forma consistente o teto shape-$t_0$.

	\paragraph{União tardia passa o teto.}
	Cinco SVMs, média de scores: forma não vê 144 GLCM; volume não vê firstorder.
	Late Q4 e a âncora superam o máximo unimodal nas quatro células.
	No claim, late só-$t_0$ empata o teto; o ganho exige os braços Q4.
	A âncora é o desenho que o unimodal pede: geometria no baseline, o resto no tempo.

	\paragraph{Fatorial, não crowning.}
	As quatro coortes não são réplicas independentes, mas o padrão (shape-$t_0$ teto; Q4 em texture/FO/vol; late $>$ teto) sobrevive a janela e intervalo.
	Não se coroa \code{wide}, deltas nem early fusion.

	\paragraph{Atributos.}
	\code{l1\_stable} reduz a classe a poucos representantes (SurfaceVolumeRatio, \code{csf\_norm}, InverseVariance, Entropy/Uniformity).
	Frequência across folds, não biomarcador único.
	Deslocamento em 12\,meses cai em \code{corr\_fallback}: DVF normativo continua um negativo útil, não um claim.

	\paragraph{Clínico.}
	Escores de $t_0$ ($0.812$) permanecem competitivos face à imagem hipocampal isolada, alinhado com~\citet{grassi-2019,moradi-2015}.
	A RM entra como adjunto; o claim das perguntas 1--3 é interno à imagem.

	\section{Limitações}
	\label{sec:limitacao}

	\begin{itemize}
		\item \textbf{Uma fonte, caso completo.} Só a ADNI fornece $n$ viável com três RM e rótulo sMCI/pMCI nas janelas 36/48 meses.
		\item \textbf{Soft pMCI.} 46 (intervalos de 6 meses) e 33 (12 meses) têm $i_3$ igual ao primeiro AD.
		\item \textbf{Sobreposição e $n$.} Coortes não independentes; \code{48m\_12m} tem 120 sujeitos.
		\item \textbf{Q4 são níveis, não taxas.} Não se modela $T_b-T_a$.
		\item \textbf{Fallback L1.} Disp em 12\,meses não é pool $\pi\ge 70\%$.
		\item \textbf{Inferência pareada.} Os $p$/$q$ destes CSVs ainda não substituem a tabela FDR do protocolo \code{t1\_deltas}; os pontos e o $\Delta$ são o resultado reportado.
		\item \textbf{ROI única.} Hipocampo de propósito.
		\item \textbf{Operações pré-CV.} QC do MRQy e templates CN estimam-se uma vez no conjunto.
	\end{itemize}

	\section{Conclusões}
	\label{sec:conclusoes}

	Sob validação aninhada no paciente, Q4 (\code{t1\_d21\_d32}) \textbf{agrega} discriminação unimodal em volume, textura e primeira ordem, mas \textbf{não} supera o teto da forma em $t_0$.
	A união tardia das cinco famílias --- em especial a âncora forma@$t_0$ $\cup$ restante Q4 --- \textbf{supera} esse teto nas quatro células do fatorial ($0.76\to 0.80$ em \code{48m\_12m}).
	\code{l1\_stable} retém poucos representantes estáveis, não uma assinatura única.
	Escores clínicos de $t_0$ permanecem competitivos; a RM hipocampal é adjunto estrutural.
	O contributo é o protocolo auditável e esta leitura por papéis, não um classificador novo.

	\section*{Agradecimentos}

	A recolha e a partilha dos dados foram financiadas pela Alzheimer's Disease Neuroimaging Initiative (ADNI).
	
	\subsection*{Financiamento}
	
	O presente trabalho foi realizado com apoio da Coordenação de Aperfeiçoamento de Pessoal de Nível Superior -- Brasil (CAPES) -- Código de Financiamento 001.
	
	\section*{Conflito de interesses}
	
	Os autores declaram não haver conflito de interesses.
	
	\bibliographystyle{cas-model2-names-short}
	\biboptions{authoryear}
	\bibliography{referencias}
\end{document}
